% Complete reviewed Unit 11; editorial master: chapters/ch11/answer-key.md
% Markdown SHA-256: fc0c8d0667c675f57673606e3d05bc3b7e6a75adc7c9a6d21e40bb747af58d42
% One-use initial assembly followed by complete content and page review.
\newcommand{\runningtitle}{Unit 11. Answer key}
% Layout inherited from the reviewed previous course; set a running title per unit.
% This is a small hand-maintained TeX layout, not a conversion pipeline.
\RequirePackage{fix-cm}
\documentclass[11pt,twoside]{article}
\usepackage{fontspec}
\setmainfont{Cambria}[
  Renderer=HarfBuzz,
  BoldFont=cambriab.ttf,
  ItalicFont=cambriai.ttf,
  BoldItalicFont=cambriaz.ttf
]
\setmonofont{lmmono10-regular.otf}[BoldFont=lmmonolt10-bold.otf]
\usepackage[paperwidth=165mm,paperheight=241mm,inner=21mm,outer=19mm,
  top=22mm,bottom=20mm,headheight=14pt,headsep=7mm,footskip=10mm]{geometry}
\usepackage{xcolor,array,booktabs,longtable,tabularx}
\usepackage{enumitem,fancyhdr}
\usepackage[unicode,hidelinks]{hyperref}
\definecolor{ink}{HTML}{000000}
\definecolor{muted}{HTML}{4C5558}
\setlength{\parindent}{0pt}
\setlength{\parskip}{3.5pt}
\AtBeginDocument{\fontsize{11}{14.3}\selectfont}
\setlength{\tabcolsep}{4pt}
\renewcommand{\arraystretch}{1.10}
\setlist[itemize]{leftmargin=1.3em,label=\textbullet,itemsep=2pt,topsep=3pt,parsep=0pt}
\setlist[enumerate]{leftmargin=1.8em,itemsep=3pt,topsep=4pt,parsep=0pt}
\widowpenalty=10000
\clubpenalty=10000
\emergencystretch=1em
\raggedbottom
\pagestyle{fancy}
\fancyhf{}
\providecommand{\runningtitle}{Old English}
\fancyhead[LE,RO]{\small\itshape \runningtitle}
\fancyhead[LO,RE]{\small Old English}
\fancyfoot[LE,RO]{\small\thepage}
\renewcommand{\headrulewidth}{0.3pt}
\newcommand{\booktitle}[2]{
  {\fontsize{9}{11}\selectfont\color{muted}#1\par}
  \vspace{3pt}
  {\fontsize{20}{24}\selectfont\color{ink}#2\par}
  \vspace{6pt}
}
\newcommand{\source}[1]{\par{\footnotesize\color{muted}#1\par}}
\newcommand{\readingtext}[1]{
  \par\addvspace{5pt}
  \begingroup\leftskip=1.5em
    \fontsize{12}{16}\selectfont #1\par
  \endgroup\addvspace{5pt}
}
\makeatletter
\renewcommand{\section}{\@startsection{section}{1}{\z@}
  {-2.4ex \@plus -0.6ex \@minus -0.2ex}{0.8ex \@plus 0.2ex}
  {\normalfont\fontsize{12}{15}\selectfont\bfseries\color{ink}}}
\renewcommand{\subsection}{\@startsection{subsection}{2}{\z@}
  {-1.8ex \@plus -0.5ex \@minus -0.2ex}{0.5ex \@plus 0.2ex}
  {\normalfont\fontsize{11}{14}\selectfont\bfseries\color{ink}}}
\makeatother

\hyphenpenalty=10000
\exhyphenpenalty=10000
\emergencystretch=3em
\newcommand{\itempara}[2]{\par\addvspace{3pt}\noindent\hangindent=1.8em\hangafter=1\makebox[1.8em][l]{#1.}#2\par}
\newcommand{\reserve}[1]{\par\penalty-100\begingroup\dimen0=\pagegoal\advance\dimen0 by-\pagetotal\ifdim\dimen0<#1\newpage\fi\endgroup}
\begin{document}
\addfontfeatures{Ligatures=TeXOff}
\raggedright
\booktitle{UNIT 11}{Answer key}

% mdparagraph:P001:begin
These answers preserve each task's meaning and stated construction. Other grammatical word orders and documented forms are acceptable. Where a task specifies a conjunction, selected paradigm or viewpoint, distinguish a valid alternative from an answer that meets that particular target. Equivalent English translations need not reproduce the model's wording.
% mdparagraph:P001:end

% mdparagraph:P002:begin
The key covers all 92 tasks. Cook's three complete sets retain their original English prompts and numbering; the Old English answers and explanations are new teaching material. The fourteen translated teaching examples are not additional exercises. [1]
% mdparagraph:P002:end

% mdparagraph:P003:begin
In complete form tasks, supply every labeled cell even when several cells share a spelling. In translations, agreement follows grammatical gender and number. The new participles are base forms: the earlier bare-predicate passive permits the uninflected bases used below, including with plural subjects. Correctly evidenced inflected predicates may also be valid, but an extra ending is not required merely because the subject is plural.
% mdparagraph:P003:end

% mdparagraph:P004:begin
Mood must be read with the clause's function and context. A purpose is an intended outcome; an asserted result is presented as occurring. A concessive subjunctive can admit a fact. A preterite subjunctive in the stated contrary-present condition refers to the present. Shared spellings, especially first-singular forms, do not identify mood by themselves. The local Cook and reading helps remain available: supplied expressions are not tests of unaided recall. [2]
% mdparagraph:P004:end

\clearpage
% mdparagraph:P005:begin
\section*{F. Complete forms}
% mdparagraph:P005:end

\reserve{9\baselineskip}
\par\begin{minipage}{\linewidth}\raggedright
% mdparagraph:P006:begin
\textbf{F01. Bæð, geat and sīðfæt.}
% mdparagraph:P006:end

% item: NEW-U11-F01
% key: KEY-NEW-U11-F01
% mdtable:T01:begin
\par\begingroup
\fontsize{10.5}{13.2}\selectfont
\setlength{\LTleft}{0pt}\setlength{\LTright}{\fill}
\noindent
\begin{tabular}{@{}>{\raggedright\arraybackslash}p{22.44mm}>{\raggedright\arraybackslash}p{22.44mm}>{\raggedright\arraybackslash}p{22.44mm}>{\raggedright\arraybackslash}p{23.0mm}>{\raggedright\arraybackslash}p{22.44mm}@{}}
\toprule
\textbf{Case} & \textbf{bæð singular} & \textbf{bæð plural} & \textbf{geat singular} & \textbf{geat plural} \\
\midrule
Nominative & bæð & baðu & geat & gatu \\
Genitive & bæðes & baða & geates & gata \\
Dative & bæðe & baðum & geate & gatum \\
Accusative & bæð & baðu & geat & gatu \\
\bottomrule
\end{tabular}
\par\endgroup

% mdtable:T01:end
% mdtable:T02:begin
\par\begingroup
\fontsize{10.5}{13.2}\selectfont
\setlength{\LTleft}{0pt}\setlength{\LTright}{\fill}
\noindent
\begin{tabular}{@{}>{\raggedright\arraybackslash}p{72.62mm}>{\raggedright\arraybackslash}p{48.57mm}@{}}
\toprule
\textbf{Requested sīðfæt cell} & \textbf{Form} \\
\midrule
Nominative plural & sīðfatas \\
Dative plural & sīðfatum \\
\bottomrule
\end{tabular}
\par\endgroup

% mdtable:T02:end
% mdparagraph:P007:begin
\textbf{Explanation:} Bæð keeps æ throughout the singular and uses bað- throughout the selected plural; geat similarly has singular geat- and plural gat-. Sīðfæt has its own specified plural stem sīðfat-. These changes are lexical, not a rule changing every æ or ea before an ending. The attested dative plural geatum is a valid alternative to selected gatum; it does not replace the other selected plural cells. Do not give masculine -as plurals to the two neuters. [3]
% mdparagraph:P007:end

\end{minipage}\par\addvspace{4pt}

\par\begin{minipage}{\linewidth}\raggedright
% mdparagraph:P008:begin
\textbf{F02. Complete seolh.}
% mdparagraph:P008:end

% item: NEW-U11-F02
% key: KEY-NEW-U11-F02
% mdtable:T03:begin
\par\begingroup
\fontsize{10.5}{13.2}\selectfont
\setlength{\LTleft}{0pt}\setlength{\LTright}{\fill}
\noindent
\begin{tabular}{@{}>{\raggedright\arraybackslash}p{38.56mm}>{\raggedright\arraybackslash}p{41.26mm}>{\raggedright\arraybackslash}p{38.56mm}@{}}
\toprule
\textbf{Case} & \textbf{seolh singular} & \textbf{seolh plural} \\
\midrule
Nominative & seolh & sēolas \\
Genitive & sēoles & sēola \\
Dative & sēole & sēolum \\
Accusative & seolh & sēolas \\
\bottomrule
\end{tabular}
\par\endgroup

% mdtable:T03:end
% mdparagraph:P009:begin
\textbf{Explanation:} Nominative and accusative singular seolh retain h and short eo. Every ending-bearing cell uses sēol-, with both h-loss and vowel lengthening. A differently ordered table is equivalent if its labels are correct. Leaving h before the endings or retaining short eo loses part of the requested alternation; do not transfer a different h-noun's pattern to this lexical system. [3]
% mdparagraph:P009:end

\end{minipage}\par\addvspace{4pt}

\par\begin{minipage}{\linewidth}\raggedright
% mdparagraph:P010:begin
\textbf{F03. Complete sumer and winter.}
% mdparagraph:P010:end

% item: NEW-U11-F03
% key: KEY-NEW-U11-F03
% mdtable:T04:begin
\par\begingroup
\fontsize{10.5}{13.2}\selectfont
\setlength{\LTleft}{0pt}\setlength{\LTright}{\fill}
\noindent
\begin{tabular}{@{}>{\raggedright\arraybackslash}p{21.75mm}>{\raggedright\arraybackslash}p{23.27mm}>{\raggedright\arraybackslash}p{21.75mm}>{\raggedright\arraybackslash}p{24.25mm}>{\raggedright\arraybackslash}p{21.75mm}@{}}
\toprule
\textbf{Case} & \textbf{sumer singular} & \textbf{sumer plural} & \textbf{winter singular} & \textbf{winter plural} \\
\midrule
Nominative & sumer & sumeras & winter & wintras \\
Genitive & sumeres & sumera & wintres & wintra \\
Dative & sumere & sumerum & wintre & wintrum \\
Accusative & sumer & sumeras & winter & wintras \\
\bottomrule
\end{tabular}
\par\endgroup

% mdtable:T04:end
% mdparagraph:P011:begin
\textbf{Explanation:} Sumer retains its medial vowel before these endings; winter loses the medial e, giving wintr-. The documented alternatives sumera and wintra may each serve genitive singular, dative singular or nominative plural. They are acceptable in those cells and need not be recalled in addition to the selected table. Their existence does not authorize -a in every case, or syncope in sumer. [3]
% mdparagraph:P011:end

\end{minipage}\par\addvspace{4pt}

\par\begin{minipage}{\linewidth}\raggedright
% mdparagraph:P012:begin
\textbf{F04. Diagnostic adjective forms.}
% mdparagraph:P012:end

% item: NEW-U11-F04
% key: KEY-NEW-U11-F04
% mdtable:T05:begin
\par\begingroup
\fontsize{10.5}{13.2}\selectfont
\setlength{\LTleft}{0pt}\setlength{\LTright}{\fill}
\noindent
\begin{tabular}{@{}>{\raggedright\arraybackslash}p{83.32mm}>{\raggedright\arraybackslash}p{37.87mm}@{}}
\toprule
\textbf{Word and requested cell} & \textbf{Form} \\
\midrule
Blæc: strong masculine accusative singular & blæcne \\
Blæc: strong feminine nominative singular & blacu \\
Blæc: weak masculine nominative singular & blaca \\
Stille: strong masculine accusative singular & stilne \\
Fægerra: masculine nominative singular & fægerra \\
Fægerra: feminine nominative singular & fægerre \\
Fægerra: neuter nominative singular & fægerre \\
Dēorwierðra: masculine dative plural & dēorwierðrum \\
Dēorwierðra: feminine dative plural & dēorwierðrum \\
Dēorwierðra: neuter dative plural & dēorwierðrum \\
\bottomrule
\end{tabular}
\par\endgroup

% mdtable:T05:end
% mdparagraph:P013:begin
\textbf{Explanation:} Blæc has æ before consonantal -ne, but a before the selected vowel endings. Stille simplifies ll before consonantal -ne: stilne. Both comparatives use weak endings; the three dative plurals coincide. The selected forms, rather than a different synonym or degree of comparison, answer this form task. Do not retain ll in stilne or substitute strong endings for the comparatives. [3]
% mdparagraph:P013:end

\end{minipage}\par\addvspace{4pt}

\par\begin{minipage}{\linewidth}\raggedright
% mdparagraph:P014:begin
\textbf{F05. Two complete indicative systems.}
% mdparagraph:P014:end

% item: NEW-U11-F05
% key: KEY-NEW-U11-F05
% mdparagraph:P015:begin
\textbf{scęððan}
% mdparagraph:P015:end

% mdtable:T06:begin
% course-table-style: 2026-10
% course-panel:T06-P01:begin
\par\begingroup
\fontsize{10.5}{13.2}\selectfont
\setlength{\LTleft}{0pt}\setlength{\LTright}{\fill}
\noindent
\begin{tabular}{@{}>{\raggedright\arraybackslash}p{39.459mm}>{\raggedright\arraybackslash}p{39.459mm}>{\raggedright\arraybackslash}p{39.459mm}@{}}
\toprule
\textbf{Person and number} & \textbf{Present} & \textbf{Preterite} \\
\midrule
1 sg. & scęððe & scęðede \\
2 sg. & scęðest & scęðedest \\
3 sg. & scęðeð & scęðede \\
1–3 pl. & scęððað & scęðedon \\
\bottomrule
\end{tabular}
\par\endgroup
% course-panel:T06-P01:end
% mdtable:T06:end
% mdparagraph:P016:begin
\textbf{genēalǣcan}
% mdparagraph:P016:end

% mdtable:T07:begin
% course-table-style: 2026-10
% course-panel:T07-P01:begin
\par\begingroup
\fontsize{10.5}{13.2}\selectfont
\setlength{\LTleft}{0pt}\setlength{\LTright}{\fill}
\noindent
\begin{tabular}{@{}>{\raggedright\arraybackslash}p{39.459mm}>{\raggedright\arraybackslash}p{39.459mm}>{\raggedright\arraybackslash}p{39.459mm}@{}}
\toprule
\textbf{Person and number} & \textbf{Present} & \textbf{Preterite} \\
\midrule
1 sg. & genēalǣce & genēalǣcte \\
2 sg. & genēalǣcst & genēalǣctest \\
3 sg. & genēalǣcð & genēalǣcte \\
1–3 pl. & genēalǣcað & genēalǣcton \\
\bottomrule
\end{tabular}
\par\endgroup
% course-panel:T07-P01:end
% mdtable:T07:end
% mdparagraph:P017:begin
\textbf{Explanation:} Scęððan uses the selected weak system: the present retains ðð in scęððe/\allowbreak{}scęððað, whereas scęðest/\allowbreak{}scęðeð and the dental preterite scęðede have single ð. The retained ending vowels distinguish this selected short-stem pattern. Genēalǣcan has present -c- and selected preterite -ct-; its documented -ht- alternative, as in genēalǣhte, is also valid. Both preterites have second-singular -est and plural -on. Do not invent a strong past for the requested scęððan table or let identical plural forms erase their three person labels. [4]
% mdparagraph:P017:end

\end{minipage}\par\addvspace{4pt}

\par\begin{minipage}{\linewidth}\raggedright
% mdparagraph:P018:begin
\textbf{F06. Feohtan and bringan diagnostics.}
% mdparagraph:P018:end

% item: NEW-U11-F06
% key: KEY-NEW-U11-F06
% mdtable:T08:begin
\par\begingroup
\fontsize{10.5}{13.2}\selectfont
\setlength{\LTleft}{0pt}\setlength{\LTright}{\fill}
\noindent
\begin{tabular}{@{}>{\raggedright\arraybackslash}p{53.95mm}>{\raggedright\arraybackslash}p{32.21mm}>{\raggedright\arraybackslash}p{32.21mm}@{}}
\toprule
\textbf{Requested indicative cell} & \textbf{Feohtan} & \textbf{Bringan} \\
\midrule
Present second singular & fiehtst & bringst \\
Present third singular & fieht & bringð \\
Preterite first singular & feaht & brōhte \\
Preterite second singular & fuhte & brōhtest \\
Preterite first plural & fuhton & brōhton \\
\bottomrule
\end{tabular}
\par\endgroup

% mdtable:T08:end
% mdparagraph:P019:begin
\textbf{Explanation:} Feohtan has a strong vowel series: first-singular feaht differs from second-singular fuhte and plural fuhton. Bringan has the dental preterite brōht-, with long ō, despite its strong-looking present. The documented full or unmutated present variants of feohtan are valid language; the displayed fiehtst/\allowbreak{}fieht are the requested selected forms. Do not make the second-singular strong past feaht or remove the dental formation from brōhte. [4]
% mdparagraph:P019:end

\end{minipage}\par\addvspace{4pt}

\par\begin{minipage}{\linewidth}\raggedright
% mdparagraph:P020:begin
\textbf{F07. Present and past participle bases.}
% mdparagraph:P020:end

% item: NEW-U11-F07
% key: KEY-NEW-U11-F07
% mdtable:T09:begin
\par\begingroup
\fontsize{10.5}{13.2}\selectfont
\setlength{\LTleft}{0pt}\setlength{\LTright}{\fill}
\noindent
\begin{tabular}{@{}>{\raggedright\arraybackslash}p{32.04mm}>{\raggedright\arraybackslash}p{47.09mm}>{\raggedright\arraybackslash}p{39.25mm}@{}}
\toprule
\textbf{Verb} & \textbf{Present participle} & \textbf{Past participle} \\
\midrule
Fōn & fōnde & fangen \\
Findan & findende & funden \\
Fēdan & fēdende & fēded \\
Geweorðian & geweorðiende & geweorðod \\
Hȳdan & hȳdende & hȳded \\
\bottomrule
\end{tabular}
\par\endgroup

% mdtable:T09:end
% mdparagraph:P021:begin
\textbf{Explanation:} The ordinary present-participle ending is -ende, with -iende in the weak-II stem geweorðiende; contracted fōn has fōnde. Strong fangen/\allowbreak{}funden have -en and their lexical past-participle stems. The weak past bases have dental -ed or -od. Contracted fēd and hȳd are defensible past-participle variants when that analysis is clear. Do not produce fōnende, confuse fangen with the finite past fēng, or substitute a finite third-person ending for a participle ending. [4]
% mdparagraph:P021:end

\end{minipage}\par\addvspace{4pt}

\par\begin{minipage}{\linewidth}\raggedright
% mdparagraph:P022:begin
\textbf{F08. Commands and number contrasts.}
% mdparagraph:P022:end

% item: NEW-U11-F08
% key: KEY-NEW-U11-F08
% mdtable:T10:begin
\par\begingroup
\fontsize{10.5}{13.2}\selectfont
\setlength{\LTleft}{0pt}\setlength{\LTright}{\fill}
\noindent
\begin{tabular}{@{}>{\raggedright\arraybackslash}p{66.23mm}>{\raggedright\arraybackslash}p{26.07mm}>{\raggedright\arraybackslash}p{26.07mm}@{}}
\toprule
\textbf{System} & \textbf{Singular} & \textbf{Plural} \\
\midrule
Gewītan: command & gewīt & gewītað \\
Hȳdan: preterite indicative third person & hȳdde & hȳddon \\
Scēawian: present indicative third person & scēawað & scēawiað \\
Geweorðian: present indicative third person & geweorðað & geweorðiað \\
\bottomrule
\end{tabular}
\par\endgroup

% mdtable:T10:end
% mdparagraph:P023:begin
\textbf{Explanation:} Gewīt addresses one person; gewītað addresses several. The past third-person contrast is -e/\allowbreak{}-on after hȳdd-. The two weak-II presents distinguish singular -að from plural -iað. The form labels matter even where a command and an indicative can share a surface. Another verb meaning go would not supply the requested gewītan cells; do not omit the i that distinguishes these weak-II plurals. [4]
% mdparagraph:P023:end

\end{minipage}\par\addvspace{4pt}

\reserve{7\baselineskip}
% mdparagraph:P024:begin
\section*{A. Clause analysis}
% mdparagraph:P024:end

\reserve{5\baselineskip}
\par\begin{minipage}{\linewidth}\raggedright
% mdparagraph:P025:begin
\textbf{A01. Past when and then.}
% mdparagraph:P025:end

% item: NEW-U11-A01
% key: KEY-NEW-U11-A01
% mdparagraph:P026:begin
\textbf{Ðā se munuc cōm, ðā sealde se mann him wæter.}
% mdparagraph:P026:end

% mdparagraph:P027:begin
When the monk came, the man gave him water.
% mdparagraph:P027:end

% mdparagraph:P028:begin
The subordinate temporal clause is \textbf{ðā se munuc cōm}; cōm is third-singular preterite indicative. The main clause is \textbf{ðā sealde se mann him wæter}; sealde is third-singular preterite indicative in this factual narration. The first ðā means when; the second correlates with it as then. English may omit then. Him is the monk, the dative recipient. Neither ðā is an article here, and the spelling sealde alone does not prove subjunctive mood.
% mdparagraph:P028:end

\end{minipage}\par\addvspace{4pt}

\par\begin{minipage}{\linewidth}\raggedright
% mdparagraph:P029:begin
\textbf{A02. Repeated visits.}
% mdparagraph:P029:end

% item: NEW-U11-A02
% key: KEY-NEW-U11-A02
% mdparagraph:P030:begin
\textbf{Ðonne gē cumað, ðonne helpe ic ēow.}
% mdparagraph:P030:end

% mdparagraph:P031:begin
Whenever you come, I help you.
% mdparagraph:P031:end

% mdparagraph:P032:begin
\textbf{Ðonne gē cumað} is the temporal subordinate clause, with second-plural present indicative cumað. The second ðonne is the main-clause then; helpe is first-singular present indicative in the stated recurring situation. Ēow is the plural dative person helped. When you come, then I help you is equivalent. Translating ðonne as than misses the temporal relation; the first-singular spelling helpe alone cannot distinguish indicative from subjunctive.
% mdparagraph:P032:end

\end{minipage}\par\addvspace{4pt}

\par\begin{minipage}{\linewidth}\raggedright
% mdparagraph:P033:begin
\textbf{A03. Simultaneous actions.}
% mdparagraph:P033:end

% item: NEW-U11-A03
% key: KEY-NEW-U11-A03
% mdparagraph:P034:begin
\textbf{(a) Ðā-hwīle-ðe hē ǣt, ic bær wæter.}
% mdparagraph:P034:end

% mdparagraph:P035:begin
\textbf{(b) Mid-ðȳ hē ǣt, ic bær wæter.}
% mdparagraph:P035:end

% mdparagraph:P036:begin
Both statements mean \textbf{While he ate, I carried water}.
% mdparagraph:P036:end

% mdparagraph:P037:begin
Ðā-hwīle-ðe and mid-ðȳ introduce the accompanying temporal circumstance. Ǣt is third-singular and bær first-singular preterite indicative. Neither sentence by itself says that the eating caused the carrying. As he was eating is an idiomatic English alternative; it does not change the Old English simple-preterite analysis. Because would assert a causal relation that these words do not themselves express.
% mdparagraph:P037:end

\end{minipage}\par\addvspace{4pt}

\par\begin{minipage}{\linewidth}\raggedright
% mdparagraph:P038:begin
\textbf{A04. Before and until an anticipated arrival.}
% mdparagraph:P038:end

% item: NEW-U11-A04
% key: KEY-NEW-U11-A04
% mdparagraph:P039:begin
\textbf{(a) Ne far, ǣr ic cume.}
% mdparagraph:P039:end

% mdparagraph:P040:begin
\textbf{(b) Wes hēr, oð-ðæt ic cume.}
% mdparagraph:P040:end

% mdparagraph:P041:begin
Do not go before I come. Be here until I come.
% mdparagraph:P041:end

% mdparagraph:P042:begin
Far and wes are singular commands. In each selected temporal clause, cume is first-singular present subjunctive: the coming is still anticipated. Its shape is identical to the first-singular present indicative, so context and construction establish this analysis. Stay here until I arrive is an equivalent English rendering. No distinct future ending or willan is required, and the spelling cume alone does not settle mood.
% mdparagraph:P042:end

\end{minipage}\par\addvspace{4pt}

\par\begin{minipage}{\linewidth}\raggedright
% mdparagraph:P043:begin
\textbf{A05. A stated reason.}
% mdparagraph:P043:end

% item: NEW-U11-A05
% key: KEY-NEW-U11-A05
% mdparagraph:P044:begin
\textbf{Ic helpe ðē, for-ðon-ðe ðū eart untrum.}
% mdparagraph:P044:end

% mdparagraph:P045:begin
I help you because you are sick.
% mdparagraph:P045:end

% mdparagraph:P046:begin
\textbf{For-ðon-ðe ðū eart untrum} gives the reason; \textbf{ic helpe ðē} is the action explained by it. Eart is second-singular present indicative. Ðē is the dative person helped. Since you are sick is equivalent when since is causal. Reversing reason and action changes the statement, and so that would turn a stated fact into an intended purpose.
% mdparagraph:P046:end

\end{minipage}\par\addvspace{4pt}

\par\begin{minipage}{\linewidth}\raggedright
% mdparagraph:P047:begin
\textbf{A06. Three functions of ðæt.}
% mdparagraph:P047:end

% item: NEW-U11-A06
% key: KEY-NEW-U11-A06
% mdparagraph:P048:begin
\textbf{(a) Ic wāt ðæt se munuc cymð.}
% mdparagraph:P048:end

% mdparagraph:P049:begin
\textbf{(b) Se munuc cymð, ðæt hē mē helpe.}
% mdparagraph:P049:end

% mdparagraph:P050:begin
\textbf{(c) Se munuc is swā gōd, ðæt hē mē hilpð.}
% mdparagraph:P050:end

% mdtable:T11:begin
\par\begingroup
\fontsize{10.5}{13.2}\selectfont
\setlength{\LTleft}{0pt}\setlength{\LTright}{\fill}
\noindent
\begin{tabular}{@{}>{\raggedright\arraybackslash}p{17.45mm}>{\raggedright\arraybackslash}p{40.39mm}>{\raggedright\arraybackslash}p{60.54mm}@{}}
\toprule
\textbf{Sentence} & \textbf{Translation} & \textbf{Function and finite verb} \\
\midrule
(a) & I know that the monk comes. & Content; cymð, third-singular present indicative. \\
(b) & The monk comes so that he may help me. & Intended purpose; helpe, third-singular present subjunctive. \\
(c) & The monk is so good that he helps me. & Asserted result; hilpð, third-singular present indicative. \\
\bottomrule
\end{tabular}
\par\endgroup

% mdtable:T11:end
% mdparagraph:P051:begin
The first clause supplies what is known; the second states why he comes; the third presents helping as the result. Comes/\allowbreak{}is coming and so that he may help/\allowbreak{}in order to help are defensible English choices. They do not make the three Old English clauses interchangeable. The spelling ðæt neither fixes the function nor requires subjunctive mood in every occurrence.
% mdparagraph:P051:end

\end{minipage}\par\addvspace{4pt}

\par\begin{minipage}{\linewidth}\raggedright
% mdparagraph:P052:begin
\textbf{A07. An outcome to prevent.}
% mdparagraph:P052:end

% item: NEW-U11-A07
% key: KEY-NEW-U11-A07
% mdparagraph:P053:begin
\textbf{Se mann hēold ðæt dēor, ðȳ-lǣs-ðe hit ðone męte nōme.}
% mdparagraph:P053:end

% mdparagraph:P054:begin
The man held the animal lest it take the food.
% mdparagraph:P054:end

% mdparagraph:P055:begin
Equivalently: he held it so that it would not take the food. \textbf{Ðȳ-lǣs-ðe hit ðone męte nōme} gives the negative purpose; nōme is third-singular preterite subjunctive. The intended prevention concerns the taking, and the clause does not assert that taking occurred. The conjunction supplies the negative purpose. Negating the main verb would instead say that the man did not hold the animal.
% mdparagraph:P055:end

\end{minipage}\par\addvspace{4pt}

\par\begin{minipage}{\linewidth}\raggedright
% mdparagraph:P056:begin
\textbf{A08. A concession that admits a fact.}
% mdparagraph:P056:end

% item: NEW-U11-A08
% key: KEY-NEW-U11-A08
% mdparagraph:P057:begin
\textbf{Ðēah hē sīe untrum, hē hilpð mē.}
% mdparagraph:P057:end

% mdparagraph:P058:begin
Although he is sick, he helps me.
% mdparagraph:P058:end

% mdparagraph:P059:begin
Sīe is third-singular present subjunctive of wesan. The concessive clause admits the illness as a circumstance despite which he helps. It is compatible with the explicitly accepted fact of his illness; subjunctive does not mean false. Though or even though gives the same relation. If would instead leave the illness as a condition and would miss this task's stated concession.
% mdparagraph:P059:end

\end{minipage}\par\addvspace{4pt}

\par\begin{minipage}{\linewidth}\raggedright
% mdparagraph:P060:begin
\textbf{A09. An open condition and an absent person.}
% mdparagraph:P060:end

% item: NEW-U11-A09
% key: KEY-NEW-U11-A09
% mdparagraph:P061:begin
\textbf{(a) Gif hē is hēr, hē hilpð mē.}
% mdparagraph:P061:end

% mdparagraph:P062:begin
\textbf{(b) Gif hē wǣre hēr, hē hulpe mē.}
% mdparagraph:P062:end

% mdtable:T12:begin
\par\begingroup
\fontsize{10.5}{13.2}\selectfont
\setlength{\LTleft}{0pt}\setlength{\LTright}{\fill}
\noindent
\begin{tabular}{@{}>{\raggedright\arraybackslash}p{19.02mm}>{\raggedright\arraybackslash}p{40.6mm}>{\raggedright\arraybackslash}p{58.76mm}@{}}
\toprule
\textbf{Sentence} & \textbf{Translation} & \textbf{Finite forms} \\
\midrule
(a) & If he is here, he helps me. & Is and hilpð: present indicative, third singular. \\
(b) & If he were here, he would help me. & Wǣre and hulpe: preterite subjunctive, third singular. \\
\bottomrule
\end{tabular}
\par\endgroup

% mdtable:T12:end
% mdparagraph:P063:begin
In (a), his presence remains open; a forward-looking he will help me can also fit its consequence. In (b), the stipulated absence is now, so both preterite subjunctives concern a contrary-present situation. They do not assert a past visit. Gif alone requires neither this subjunctive pair nor the indicative pair; English would does not compel an added willan.
% mdparagraph:P063:end

\end{minipage}\par\addvspace{4pt}

\par\begin{minipage}{\linewidth}\raggedright
% mdparagraph:P064:begin
\textbf{A10. Question, temporal clause and comparison.}
% mdparagraph:P064:end

% item: NEW-U11-A10
% key: KEY-NEW-U11-A10
% mdparagraph:P065:begin
\textbf{(a) Hwonne cymð se munuc?}
% mdparagraph:P065:end

% mdparagraph:P066:begin
\textbf{(b) Ðonne se munuc cymð, ic helpe him.}
% mdparagraph:P066:end

% mdparagraph:P067:begin
\textbf{(c) Se munuc is ieldra ðonne se fiscere.}
% mdparagraph:P067:end

% mdparagraph:P068:begin
When does the monk come? When the monk comes, I help him. The monk is older than the fisherman.
% mdparagraph:P068:end

% mdparagraph:P069:begin
Hwonne asks a direct question. The first ðonne introduces a subordinate temporal clause; the last ðonne introduces the comparison standard. In that comparison se fiscere is nominative, with the repeated copula understood. When is the monk coming? and whenever the monk comes are defensible contextual renderings of the first two. A single when or than gloss cannot account for all three constructions.
% mdparagraph:P069:end

\end{minipage}\par\addvspace{4pt}

\par\begin{minipage}{\linewidth}\raggedright
% mdparagraph:P070:begin
\textbf{A11. Who performs the intended action?}
% mdparagraph:P070:end

% item: NEW-U11-A11
% key: KEY-NEW-U11-A11
% mdparagraph:P071:begin
\textbf{(a) Se fiscere cōm tō helpanne ðǣm męnn.}
% mdparagraph:P071:end

% mdparagraph:P072:begin
\textbf{(b) Se fiscere cōm, ðæt se munuc ðǣm męnn hulpe.}
% mdparagraph:P072:end

% mdparagraph:P073:begin
The fisherman came to help the man. The fisherman came so that the monk might help the man.
% mdparagraph:P073:end

% mdparagraph:P074:begin
In the first sentence the understood helper in \textbf{tō helpanne} is the fisherman, the subject of cōm. In the second, \textbf{se munuc} is the explicit subject of finite hulpe, third-singular preterite subjunctive. \textbf{Ðǣm męnn} is the dative person helped in both. In order to and in order that are equivalent English renderings. Making the fisherman the helper in the second sentence would ignore its expressed purpose-clause subject.
% mdparagraph:P074:end

\end{minipage}\par\addvspace{4pt}

\par\begin{minipage}{\linewidth}\raggedright
% mdparagraph:P075:begin
\textbf{A12. The day when something happened.}
% mdparagraph:P075:end

% item: NEW-U11-A12
% key: KEY-NEW-U11-A12
% mdparagraph:P076:begin
\textbf{Hit is se dæg ðe hē cōm.}
% mdparagraph:P076:end

% mdparagraph:P077:begin
It is the day when he came.
% mdparagraph:P077:end

% mdparagraph:P078:begin
Equivalently: it is the day on which he came. \textbf{Se dæg} is the temporal antecedent of relative ðe; cōm is third-singular preterite indicative. The clause identifies the day by the event occurring on it. It asks no question and sets no condition. Replacing ðe with interrogative hwonne merely because both can be translated when would lose the relative construction.
% mdparagraph:P078:end

\end{minipage}\par\addvspace{4pt}

\reserve{7\baselineskip}
% mdparagraph:P079:begin
\section*{T. Controlled composition}
% mdparagraph:P079:end

\reserve{5\baselineskip}
\par\begin{minipage}{\linewidth}\raggedright
% mdparagraph:P080:begin
\textbf{T01. When the father came.}
% mdparagraph:P080:end

% item: NEW-U11-T01
% key: KEY-NEW-U11-T01
% mdparagraph:P081:begin
Translate using past ðā … ðā: When the father came, the son gave him water.
% mdparagraph:P081:end

% mdparagraph:P082:begin
\textbf{Ðā se fæder cōm, ðā sealde se sunu him wæter.}
% mdparagraph:P082:end

% mdparagraph:P083:begin
Cōm is third-singular preterite; sealde is the factual past giving. The father is the first clause's subject and the dative recipient him; the son gives. \textbf{Ðǣm fæder} can replace him, and grammatical clause orders may differ, but both requested ðā words must remain. Reversing giver and recipient changes the event.
% mdparagraph:P083:end

\end{minipage}\par\addvspace{4pt}

\par\begin{minipage}{\linewidth}\raggedright
% mdparagraph:P084:begin
\textbf{T02. A recurring welcome.}
% mdparagraph:P084:end

% item: NEW-U11-T02
% key: KEY-NEW-U11-T02
% mdparagraph:P085:begin
Translate as a recurring event, using ðonne … ðonne: When the fishermen come, the king gives them food.
% mdparagraph:P085:end

% mdparagraph:P086:begin
\textbf{Ðonne ðā fisceras cumað, ðonne sęlð se cyning him męte.}
% mdparagraph:P086:end

% mdparagraph:P087:begin
Cumað agrees with plural fisceras; sęlð agrees with singular cyning. Him is dative plural, the fishermen receiving food. An explicit \textbf{ðǣm fiscerum} is equivalent. The specified ðonne … ðonne presents a recurring relation; past-event ðā or singular cymð would fail the requested meaning or agreement.
% mdparagraph:P087:end

\end{minipage}\par\addvspace{4pt}

\par\begin{minipage}{\linewidth}\raggedright
% mdparagraph:P088:begin
\textbf{T03. While the fishermen sat.}
% mdparagraph:P088:end

% item: NEW-U11-T03
% key: KEY-NEW-U11-T03
% mdparagraph:P089:begin
Translate with mid-ðȳ: While the fishermen sat in the house, the son carried water.
% mdparagraph:P089:end

% mdparagraph:P090:begin
\textbf{Mid-ðȳ ðā fisceras sǣton on ðǣm hūse, se sunu bær wæter.}
% mdparagraph:P090:end

% mdparagraph:P091:begin
Sǣton is third-plural preterite; bær is third-singular preterite. On ðǣm hūse is the stationary location, with dative hūse. A correlating ðā can begin the main clause. Ðā-hwīle-ðe could express the same relation, but the task expressly requests mid-ðȳ. The carrying is simultaneous with the sitting, not its intended purpose.
% mdparagraph:P091:end

\end{minipage}\par\addvspace{4pt}

\par\begin{minipage}{\linewidth}\raggedright
% mdparagraph:P092:begin
\textbf{T04. Before I come.}
% mdparagraph:P092:end

% item: NEW-U11-T04
% key: KEY-NEW-U11-T04
% mdparagraph:P093:begin
Address one person. The speaker has not yet arrived. Translate with ǣr: Do not go before I come.
% mdparagraph:P093:end

% mdparagraph:P094:begin
\textbf{Ne far, ǣr ic cume.}
% mdparagraph:P094:end

% mdparagraph:P095:begin
Far is the command to one person. Cume is first-singular present subjunctive in this selected anticipated-before clause, although its surface also occurs in the indicative. A different grammatical order is acceptable; plural farað would address the wrong number. The arrival is expected, so an English future setting does not require an extra future auxiliary.
% mdparagraph:P095:end

\end{minipage}\par\addvspace{4pt}

\par\begin{minipage}{\linewidth}\raggedright
% mdparagraph:P096:begin
\textbf{T05. Until the fishermen come.}
% mdparagraph:P096:end

% item: NEW-U11-T05
% key: KEY-NEW-U11-T05
% mdparagraph:P097:begin
Address several people. The fishermen have not yet arrived. Translate with oð-ðæt: Be in the house until the fishermen come.
% mdparagraph:P097:end

% mdparagraph:P098:begin
\textbf{Wesað on ðǣm hūse, oð-ðæt ðā fisceras cumen.}
% mdparagraph:P098:end

% mdparagraph:P099:begin
Wesað addresses several people; \textbf{bēoð} is also an available plural command. On ðǣm hūse has the locative dative. Cumen is third-plural present subjunctive in the selected anticipated-until clause; ordinary indicative cumað would not express the specified mood target. Singular wes or cume would lose a required number distinction.
% mdparagraph:P099:end

\end{minipage}\par\addvspace{4pt}

\par\begin{minipage}{\linewidth}\raggedright
% mdparagraph:P100:begin
\textbf{T06. Because the man is thirsty.}
% mdparagraph:P100:end

% item: NEW-U11-T06
% key: KEY-NEW-U11-T06
% mdparagraph:P101:begin
Translate with for-ðon-ðe and the adjective ðurstig: The monk gives the man water because the man is thirsty.
% mdparagraph:P101:end

% mdparagraph:P102:begin
\textbf{Se munuc sęlð ðǣm męnn wæter, for-ðon-ðe se mann is ðurstig.}
% mdparagraph:P102:end

% mdparagraph:P103:begin
The man receiving water is dative \textbf{ðǣm męnn}; the causal clause has nominative \textbf{se mann} and indicative is. Masculine singular ðurstig is the agreeing predicate. A clear hē may replace the repeated se mann, and clause order may vary. Do not make the person receiving water accusative or use a purpose subjunctive for this stated reason.
% mdparagraph:P103:end

\end{minipage}\par\addvspace{4pt}

\par\begin{minipage}{\linewidth}\raggedright
% mdparagraph:P104:begin
\textbf{T07. A different subject in the purpose clause.}
% mdparagraph:P104:end

% item: NEW-U11-T07
% key: KEY-NEW-U11-T07
% mdparagraph:P105:begin
Translate with a finite ðæt purpose clause: The fishermen came so that the monk might help them.
% mdparagraph:P105:end

% mdparagraph:P106:begin
\textbf{Ðā fisceras cōmon, ðæt se munuc him hulpe.}
% mdparagraph:P106:end

% mdparagraph:P107:begin
Cōmon has the plural fishermen as subject; hulpe is third-singular preterite subjunctive with se munuc. Him is dative plural, the people helped; \textbf{ðǣm fiscerum} makes the recipient explicit. An ordinary same-subject purpose infinitive would make the fishermen the helpers and would miss the required finite clause with a different subject.
% mdparagraph:P107:end

\end{minipage}\par\addvspace{4pt}

\par\begin{minipage}{\linewidth}\raggedright
% mdparagraph:P108:begin
\textbf{T08. An actual result.}
% mdparagraph:P108:end

% item: NEW-U11-T08
% key: KEY-NEW-U11-T08
% mdparagraph:P109:begin
Translate the result as an actual past event, using swā … ðæt: The men were so hungry that they ate the food.
% mdparagraph:P109:end

% mdparagraph:P110:begin
\textbf{Ðā męnn wǣron swā hungrige, ðæt hīe ǣton ðone męte.}
% mdparagraph:P110:end

% mdparagraph:P111:begin
Hungrige is strong masculine nominative plural, agreeing with męnn. The eating is asserted to have happened: \textbf{ǣton} is third-plural preterite indicative. Hīe may be replaced by a clear repeated noun phrase. A purpose interpretation would change the given result, and singular hungrig would lose the selected adjective agreement.
% mdparagraph:P111:end

\end{minipage}\par\addvspace{4pt}

\par\begin{minipage}{\linewidth}\raggedright
% mdparagraph:P112:begin
\textbf{T09. Preventing the animals from taking food.}
% mdparagraph:P112:end

% item: NEW-U11-T09
% key: KEY-NEW-U11-T09
% mdparagraph:P113:begin
Translate with ðȳ-lǣs-ðe: The fishermen held the animals, lest the animals take our food.
% mdparagraph:P113:end

% mdparagraph:P114:begin
\textbf{Ðā fisceras hēoldon ðā dēor, ðȳ-lǣs-ðe ðā dēor ūrne męte nōmen.}
% mdparagraph:P114:end

% mdparagraph:P115:begin
Hēoldon is third-plural preterite indicative; nōmen is third-plural preterite subjunctive in the negative purpose. Ūrne agrees with singular masculine accusative męte, not with the animals. A clear hīe may replace the second ðā dēor. The conjunction means lest/\allowbreak{}so that not; the sentence does not assert that the animals actually took the food.
% mdparagraph:P115:end

\end{minipage}\par\addvspace{4pt}

\par\begin{minipage}{\linewidth}\raggedright
% mdparagraph:P116:begin
\textbf{T10. Illness admitted in a concession.}
% mdparagraph:P116:end

% item: NEW-U11-T10
% key: KEY-NEW-U11-T10
% mdparagraph:P117:begin
The monk was in fact sick. Translate using the selected concessive pattern: Although the monk was sick, he helped the man.
% mdparagraph:P117:end

% mdparagraph:P118:begin
\textbf{Ðēah se munuc wǣre untrum, hē healp ðǣm męnn.}
% mdparagraph:P118:end

% mdparagraph:P119:begin
Wǣre is third-singular preterite subjunctive in the concession; healp is third-singular preterite indicative in the stated event. \textbf{Ðēah ðe} is a documented alternative to ðēah. The person helped is dative ðǣm męnn. The subjunctive does not deny the factual illness; replacing healp with hulpe would change the main clause's presentation.
% mdparagraph:P119:end

\end{minipage}\par\addvspace{4pt}

\par\begin{minipage}{\linewidth}\raggedright
% mdparagraph:P120:begin
\textbf{T11. The king is absent now.}
% mdparagraph:P120:end

% item: NEW-U11-T11
% key: KEY-NEW-U11-T11
% mdparagraph:P121:begin
The king is absent now. Translate the contrary-present condition: If the king were here, he would help us.
% mdparagraph:P121:end

% mdparagraph:P122:begin
\textbf{Gif se cyning wǣre hēr, hē hulpe ūs.}
% mdparagraph:P122:end

% mdparagraph:P123:begin
Both wǣre and hulpe are third-singular preterite subjunctives, referring here to the stated present situation. Ūs is the dative plural person helped. A grammatical reversal of clause order is acceptable. A factual past rendering or an added willan based solely on English would would miss the selected contrary-present construction.
% mdparagraph:P123:end

\end{minipage}\par\addvspace{4pt}

\reserve{7\baselineskip}
% mdparagraph:P124:begin
\section*{Cook, Group IV, Exercise 16}
% mdparagraph:P124:end

\reserve{5\baselineskip}
\par\begin{minipage}{\linewidth}\raggedright
% mdparagraph:P125:begin
\textbf{16.1.}
% mdparagraph:P125:end

% item: EX-IV-16-01
% key: KEY-EX-IV-16-01
% mdparagraph:P126:begin
Gold is more precious (64) than silver.
% mdparagraph:P126:end

% mdparagraph:P127:begin
\textbf{Gold is dēorwierðre ðonne seolfor.}
% mdparagraph:P127:end

% mdparagraph:P128:begin
Dēorwierðre is the weak neuter nominative singular comparative agreeing with mass noun gold. Seolfor is the parallel nominative standard after ðonne. A grammatical comparison using the taught genitive standard is also defensible. Positive dēorwierðe alone would omit more; a superlative would change the comparison.
% mdparagraph:P128:end

\end{minipage}\par\addvspace{4pt}

\par\begin{minipage}{\linewidth}\raggedright
% mdparagraph:P129:begin
\textbf{16.2.}
% mdparagraph:P129:end

% item: EX-IV-16-02
% key: KEY-EX-IV-16-02
% mdparagraph:P130:begin
Righteousness is more precious than gold.
% mdparagraph:P130:end

% mdparagraph:P131:begin
\textbf{Rihtwīsnes is dēorwierðre ðonne gold.}
% mdparagraph:P131:end

% mdparagraph:P132:begin
Dēorwierðre is weak feminine nominative singular with rihtwīsnes. Its identical surface to 16.1 does not make the subject neuter. The standard gold has its own gender; it does not control the predicate adjective. A taught genitive comparison is an alternative, but more must remain comparative.
% mdparagraph:P132:end

\end{minipage}\par\addvspace{4pt}

\par\begin{minipage}{\linewidth}\raggedright
% mdparagraph:P133:begin
\textbf{16.3.}
% mdparagraph:P133:end

% item: EX-IV-16-03
% key: KEY-EX-IV-16-03
% mdparagraph:P134:begin
The bath is hot.
% mdparagraph:P134:end

% mdparagraph:P135:begin
\textbf{Ðæt bæð is hāt.}
% mdparagraph:P135:end

% mdparagraph:P136:begin
Bæð and its article are neuter nominative singular; hāt is the agreeing strong predicate. Predicate-first Hāt is ðæt bæð preserves the meaning. Masculine se or plural baðu would fail the subject's gender or number.
% mdparagraph:P136:end

\end{minipage}\par\addvspace{4pt}

\par\begin{minipage}{\linewidth}\raggedright
% mdparagraph:P137:begin
\textbf{16.4.}
% mdparagraph:P137:end

% item: EX-IV-16-04
% key: KEY-EX-IV-16-04
% mdparagraph:P138:begin
The baths are hot.
% mdparagraph:P138:end

% mdparagraph:P139:begin
\textbf{Ðā baðu sindon hāt.}
% mdparagraph:P139:end

% mdparagraph:P140:begin
Baðu is the selected neuter plural, so the copula is plural sindon. Long-stem hāt has an endingless strong neuter nominative plural. A grammatical predicate-first order is equivalent. The absence of an adjective ending does not make hāt singular; do not add masculine plural -e by analogy.
% mdparagraph:P140:end

\end{minipage}\par\addvspace{4pt}

\par\begin{minipage}{\linewidth}\raggedright
% mdparagraph:P141:begin
\textbf{16.5.}
% mdparagraph:P141:end

% item: EX-IV-16-05
% key: KEY-EX-IV-16-05
% mdparagraph:P142:begin
The animals are living.
% mdparagraph:P142:end

% mdparagraph:P143:begin
\textbf{Ðā dēor sindon cwicu.}
% mdparagraph:P143:end

% mdparagraph:P144:begin
Cwicu means alive/\allowbreak{}living and agrees with neuter plural dēor. English are living permits this ordinary adjective predicate; it does not require a new progressive construction. Dēor is unchanged in nominative singular and plural, but the article, copula and predicate establish the plural here. Do not replace cwicu with a singular predicate.
% mdparagraph:P144:end

\end{minipage}\par\addvspace{4pt}

\par\begin{minipage}{\linewidth}\raggedright
% mdparagraph:P145:begin
\textbf{16.6.}
% mdparagraph:P145:end

% item: EX-IV-16-06
% key: KEY-EX-IV-16-06
% mdparagraph:P146:begin
Some pearls are black.
% mdparagraph:P146:end

% mdparagraph:P147:begin
\textbf{Suma męregrēotan sindon blaca.}
% mdparagraph:P147:end

% mdparagraph:P148:begin
Męregrēotan is nominative plural of the weak feminine pearl noun. Suma and blaca are the selected strong feminine plural forms; blaca uses a before the vowel ending. A different grammatical order is acceptable. Masculine plural blace, used of men in 17.7, would give the wrong agreement here.
% mdparagraph:P148:end

\end{minipage}\par\addvspace{4pt}

\par\begin{minipage}{\linewidth}\raggedright
% mdparagraph:P149:begin
\textbf{16.7.}
% mdparagraph:P149:end

% item: EX-IV-16-07
% key: KEY-EX-IV-16-07
% mdparagraph:P150:begin
The island is beautiful.
% mdparagraph:P150:end

% mdparagraph:P151:begin
\textbf{Ðæt īgland is fæger.}
% mdparagraph:P151:end

% mdparagraph:P152:begin
Īgland is neuter nominative singular and fæger its strong predicate. Fæger is ðæt īgland preserves the meaning with predicate first. The English ending of island offers no guide to Old English gender; masculine se would fail the selected noun agreement.
% mdparagraph:P152:end

\end{minipage}\par\addvspace{4pt}

\par\begin{minipage}{\linewidth}\raggedright
% mdparagraph:P153:begin
\textbf{16.8.}
% mdparagraph:P153:end

% item: EX-IV-16-08
% key: KEY-EX-IV-16-08
% mdparagraph:P154:begin
Dolphins swim through the water.
% mdparagraph:P154:end

% mdparagraph:P155:begin
\textbf{Męreswīn swimmað þurh ðæt wæter.}
% mdparagraph:P155:end

% mdparagraph:P156:begin
Męreswīn is an endingless nominative plural, and swimmað is plural. Through is þurh with accusative ðæt wæter. No definite article is required before the indefinite dolphins. Equivalent word order is acceptable; singular swimð loses the plural, while a dative path would miss þurh's selected government.
% mdparagraph:P156:end

\end{minipage}\par\addvspace{4pt}

\par\begin{minipage}{\linewidth}\raggedright
% mdparagraph:P157:begin
\textbf{16.9.}
% mdparagraph:P157:end

% item: EX-IV-16-09
% key: KEY-EX-IV-16-09
% mdparagraph:P158:begin
Here are fishes often caught.
% mdparagraph:P158:end

% mdparagraph:P159:begin
\textbf{Hēr sindon fiscas oft fangen.}
% mdparagraph:P159:end

% mdparagraph:P160:begin
Plural fiscas is the subject of the present passive; fangen is the new past-participle base used under the taught bare-predicate option. Hēr and oft retain here and often. An evidenced inflected predicate can be valid, but an added plural ending is not required. Fēngon would instead be active caught, and fōnde is the present, not past, participle.
% mdparagraph:P160:end

\end{minipage}\par\addvspace{4pt}

\par\begin{minipage}{\linewidth}\raggedright
% mdparagraph:P161:begin
\textbf{16.10.}
% mdparagraph:P161:end

% item: EX-IV-16-10
% key: KEY-EX-IV-16-10
% mdparagraph:P162:begin
The night is more beautiful than the day.
% mdparagraph:P162:end

% mdparagraph:P163:begin
\textbf{Sēo niht is fægerre ðonne se dæg.}
% mdparagraph:P163:end

% mdparagraph:P164:begin
Fægerre is weak feminine nominative singular, agreeing with niht. Se dæg is the nominative comparison standard. A taught genitive-standard comparison is also defensible. The positive fæger or a predicate governed by day's masculine gender would miss the requested comparison or agreement.
% mdparagraph:P164:end

\end{minipage}\par\addvspace{4pt}

\par\begin{minipage}{\linewidth}\raggedright
% mdparagraph:P165:begin
\textbf{16.11.}
% mdparagraph:P165:end

% item: EX-IV-16-11
% key: KEY-EX-IV-16-11
% mdparagraph:P166:begin
The city was distinguished with walls and gates.
% mdparagraph:P166:end

% mdparagraph:P167:begin
\textbf{Sēo ceaster wæs geweorðod mid weallum and gatum.}
% mdparagraph:P167:end

% mdparagraph:P168:begin
Wæs keeps Cook's past time; geweorðod means distinguished/\allowbreak{}adorned. Mid takes the dative plural equipment weallum and gatum. The attested geatum is a valid dative-plural alternative. Walls and gates do not form a compound subject, and nominative gatu would lose the case required after mid.
% mdparagraph:P168:end

\end{minipage}\par\addvspace{4pt}

\par\begin{minipage}{\linewidth}\raggedright
% mdparagraph:P169:begin
\textbf{16.12.}
% mdparagraph:P169:end

% item: EX-IV-16-12
% key: KEY-EX-IV-16-12
% mdparagraph:P170:begin
Summer comes, and then winter.
% mdparagraph:P170:end

% mdparagraph:P171:begin
\textbf{Sumer cymð, and ðonne cymð winter.}
% mdparagraph:P171:end

% mdparagraph:P172:begin
Sumer and winter are separate singular subjects. Repeating cymð supplies the second predicate left understood in English; \textbf{Sumer cymð, and ðonne winter} preserves the same ellipsis. Here ðonne is sequential then, not a subordinating when or comparative than. No extra season-noun ending is required in the nominatives.
% mdparagraph:P172:end

\end{minipage}\par\addvspace{4pt}

\par\begin{minipage}{\linewidth}\raggedright
% mdparagraph:P173:begin
\textbf{16.13.}
% mdparagraph:P173:end

% item: EX-IV-16-13
% key: KEY-EX-IV-16-13
% mdparagraph:P174:begin
The sun shines over all.
% mdparagraph:P174:end

% mdparagraph:P175:begin
\textbf{Sēo sunne scīnð ofer eall.}
% mdparagraph:P175:end

% mdparagraph:P176:begin
The model takes eall as neuter accusative singular everything after ofer. \textbf{Ofer ealle}, with masculine accusative plural all people, is also defensible if that interpretation is stated. Sunne is a singular weak feminine subject. Do not make eall an unexplained plural or let the number of the things illuminated control scīnð.
% mdparagraph:P176:end

\end{minipage}\par\addvspace{4pt}

\par\begin{minipage}{\linewidth}\raggedright
% mdparagraph:P177:begin
\textbf{16.14.}
% mdparagraph:P177:end

% item: EX-IV-16-14
% key: KEY-EX-IV-16-14
% mdparagraph:P178:begin
The old man came to the house.
% mdparagraph:P178:end

% mdparagraph:P179:begin
\textbf{Se ealda mann cōm tō ðǣm hūse.}
% mdparagraph:P179:end

% mdparagraph:P180:begin
Se ealda mann is a definite nominative phrase, hence weak ealda. Cōm is third-singular preterite; tō takes dative ðǣm hūse as the destination. Grammatical word-order changes are acceptable. Accusative hūs after tō or strong eald inside this selected definite phrase would lose a target distinction.
% mdparagraph:P180:end

\end{minipage}\par\addvspace{4pt}

\par\begin{minipage}{\linewidth}\raggedright
% mdparagraph:P181:begin
\textbf{16.15.}
% mdparagraph:P181:end

% item: EX-IV-16-15
% key: KEY-EX-IV-16-15
% mdparagraph:P182:begin
Here are also found lead and iron.
% mdparagraph:P182:end

% mdparagraph:P183:begin
\textbf{Hēr sindon ēac lēad and īsern funden.}
% mdparagraph:P183:end

% mdparagraph:P184:begin
Lēad and īsern are two coordinated nominative metal nouns; together they take plural sindon. Funden is the past-participle base, with the bare-predicate option. A documented inflected predicate is a valid alternative, but the model needs no invented agreement ending. Preserve both ēac and hēr; finite fundon would make the statement active.
% mdparagraph:P184:end

\end{minipage}\par\addvspace{4pt}

\par\begin{minipage}{\linewidth}\raggedright
% mdparagraph:P185:begin
\textbf{16.16.}
% mdparagraph:P185:end

% item: EX-IV-16-16
% key: KEY-EX-IV-16-16
% mdparagraph:P186:begin
Put not your hand into the fire.
% mdparagraph:P186:end

% mdparagraph:P187:begin
\textbf{Ne dō ðīne hand on ðæt fȳr.}
% mdparagraph:P187:end

% mdparagraph:P188:begin
The model addresses one person: negative command ne dō, accusative hand with feminine accusative ðīne, and the goal on ðæt fȳr. A consistent plural-address version, \textbf{Ne dōð ēowre hand on ðæt fȳr}, is also possible. The goal is accusative; a locative dative would express being in the fire rather than the required placement into it.
% mdparagraph:P188:end

\end{minipage}\par\addvspace{4pt}

\reserve{7\baselineskip}
% mdparagraph:P189:begin
\section*{Cook, Group IV, Exercise 17}
% mdparagraph:P189:end

\reserve{5\baselineskip}
\par\begin{minipage}{\linewidth}\raggedright
% mdparagraph:P190:begin
\textbf{17.1.}
% mdparagraph:P190:end

% item: EX-IV-17-01
% key: KEY-EX-IV-17-01
% mdparagraph:P191:begin
The soldier takes the gem.
% mdparagraph:P191:end

% mdparagraph:P192:begin
\textbf{Se cęmpa nimð ðone gimm.}
% mdparagraph:P192:end

% mdparagraph:P193:begin
Cęmpa is a weak masculine nominative singular, although its nominative ends in -a. Nimð agrees with it; ðone gimm is the accusative thing taken. Another grammatical order is acceptable. Making gimm nominative with se would fail to distinguish the patient from the soldier.
% mdparagraph:P193:end

\end{minipage}\par\addvspace{4pt}

\par\begin{minipage}{\linewidth}\raggedright
% mdparagraph:P194:begin
\textbf{17.2.}
% mdparagraph:P194:end

% item: EX-IV-17-02
% key: KEY-EX-IV-17-02
% mdparagraph:P195:begin
The rain injures the tower.
% mdparagraph:P195:end

% mdparagraph:P196:begin
\textbf{Se regn scęðeð ðǣm torre.}
% mdparagraph:P196:end

% mdparagraph:P197:begin
Scęðeð is the selected third-singular present of scęððan. The harmed tower is dative \textbf{ðǣm torre}, despite English injures taking a direct object. Equivalent ordering does not alter that government. Accusative ðone torr would miss the selected lexical frame; rain remains singular regn in its ordinary weather use.
% mdparagraph:P197:end

\end{minipage}\par\addvspace{4pt}

\par\begin{minipage}{\linewidth}\raggedright
% mdparagraph:P198:begin
\textbf{17.3.}
% mdparagraph:P198:end

% item: EX-IV-17-03
% key: KEY-EX-IV-17-03
% mdparagraph:P199:begin
The whale swims.
% mdparagraph:P199:end

% mdparagraph:P200:begin
\textbf{Se hwæl swimð.}
% mdparagraph:P200:end

% mdparagraph:P201:begin
Hwæl is masculine nominative singular, with third-singular swimð. The verb is intransitive here. Reordering the subject and verb without changing the declarative meaning can be grammatical. Plural swimmað would not agree with this singular whale.
% mdparagraph:P201:end

\end{minipage}\par\addvspace{4pt}

\par\begin{minipage}{\linewidth}\raggedright
% mdparagraph:P202:begin
\textbf{17.4.}
% mdparagraph:P202:end

% item: EX-IV-17-04
% key: KEY-EX-IV-17-04
% mdparagraph:P203:begin
Fishermen caught seals.
% mdparagraph:P203:end

% mdparagraph:P204:begin
\textbf{Fisceras fēngon sēolas.}
% mdparagraph:P204:end

% mdparagraph:P205:begin
Fisceras is the plural subject of preterite fēngon; sēolas is the accusative plural object of catching. Sēol- loses h and has long ēo before the plural ending. No definite articles are required by Cook's indefinite nouns. Retaining seolh as the plural object or using present fōð would lose the noun or tense target.
% mdparagraph:P205:end

\end{minipage}\par\addvspace{4pt}

\par\begin{minipage}{\linewidth}\raggedright
% mdparagraph:P206:begin
\textbf{17.5.}
% mdparagraph:P206:end

% item: EX-IV-17-05
% key: KEY-EX-IV-17-05
% mdparagraph:P207:begin
The hungry sheep are not fed.
% mdparagraph:P207:end

% mdparagraph:P208:begin
\textbf{Ðā hungrigan scēap ne sindon fēded.}
% mdparagraph:P208:end

% mdparagraph:P209:begin
The definite adjective is weak plural hungrigan, agreeing with neuter plural scēap. The supplied ne sindon negates the plural passive; fēded is its past-participle base. Contracted fēd is a documented participial alternative when its analysis is clear. This is a statement, not a command not to feed; English hungry sheep does not license singular agreement.
% mdparagraph:P209:end

\end{minipage}\par\addvspace{4pt}

\par\begin{minipage}{\linewidth}\raggedright
% mdparagraph:P210:begin
\textbf{17.6.}
% mdparagraph:P210:end

% item: EX-IV-17-06
% key: KEY-EX-IV-17-06
% mdparagraph:P211:begin
The man’s wife is beautiful.
% mdparagraph:P211:end

% mdparagraph:P212:begin
\textbf{Ðæs mannes wīf is fæger.}
% mdparagraph:P212:end

% mdparagraph:P213:begin
Ðæs mannes is the genitive possessor, and wīf is the neuter nominative subject. Fæger therefore has neuter singular predicate agreement. A grammatical placement of the genitive after its head can preserve the sense. Feminine agreement based on the wife's natural sex would ignore the noun's grammatical gender.
% mdparagraph:P213:end

\end{minipage}\par\addvspace{4pt}

\par\begin{minipage}{\linewidth}\raggedright
% mdparagraph:P214:begin
\textbf{17.7.}
% mdparagraph:P214:end

% item: EX-IV-17-07
% key: KEY-EX-IV-17-07
% mdparagraph:P215:begin
Some men are black.
% mdparagraph:P215:end

% mdparagraph:P216:begin
\textbf{Sume męnn sindon blace.}
% mdparagraph:P216:end

% mdparagraph:P217:begin
Sume męnn is an indefinite masculine nominative plural phrase, with strong masculine plural blace. Compare feminine blaca in 16.6. Other grammatical orders are acceptable. The singular stem mann and feminine adjective ending would each lose a required plural or agreement distinction.
% mdparagraph:P217:end

\end{minipage}\par\addvspace{4pt}

\par\begin{minipage}{\linewidth}\raggedright
% mdparagraph:P218:begin
\textbf{17.8.}
% mdparagraph:P218:end

% item: EX-IV-17-08
% key: KEY-EX-IV-17-08
% mdparagraph:P219:begin
The carpenter will clothe his son.
% mdparagraph:P219:end

% mdparagraph:P220:begin
\textbf{Se trēowwyrhta scrȳt his sunu.}
% mdparagraph:P220:end

% mdparagraph:P221:begin
The present scrȳt can carry the future reference stated in the prompt. His refers to the carpenter and modifies accusative sunu, the person clothed. An explicitly volitional translation can be valid if correctly constructed, but a form of willan is not required merely by English will. Do not substitute dative suna for the selected direct object or change the possessor to a different person.
% mdparagraph:P221:end

\end{minipage}\par\addvspace{4pt}

\par\begin{minipage}{\linewidth}\raggedright
% mdparagraph:P222:begin
\textbf{17.9.}
% mdparagraph:P222:end

% item: EX-IV-17-09
% key: KEY-EX-IV-17-09
% mdparagraph:P223:begin
Depart from me, for I am (om. a) sinful man, (om. O) Lord.
% mdparagraph:P223:end

% mdparagraph:P224:begin
\textbf{Gewīt fram mē, for-ðon-ðe ic eom synfull mann, Dryhten.}
% mdparagraph:P224:end

% mdparagraph:P225:begin
Gewīt is the singular departure command; fram mē has the dative. For-ðon-ðe introduces the stated reason with indicative eom. Bare synfull mann and direct-address Dryhten preserve Cook's instructions to omit a and O. The earlier causal for-ðām is a defensible alternative. Plural gewītað would address the wrong number in this selected Lord-address reading.
% mdparagraph:P225:end

\end{minipage}\par\addvspace{4pt}

\par\begin{minipage}{\linewidth}\raggedright
% mdparagraph:P226:begin
\textbf{17.10.}
% mdparagraph:P226:end

% item: EX-IV-17-10
% key: KEY-EX-IV-17-10
% mdparagraph:P227:begin
They shall drink, but ye shall be thirsty.
% mdparagraph:P227:end

% mdparagraph:P228:begin
\textbf{Hīe drincað, ac gē bēoð ðurstige.}
% mdparagraph:P228:end

% mdparagraph:P229:begin
Hīe and gē require different persons but both take plural verbs. The presents express the indicated future setting; bēoð is particularly natural for the prospective state. Ðurstige agrees with the addressed plural. A future-context sindon can also be defensible. Do not infer that both English shall words demand the same modal auxiliary, or leave the predicate singular.
% mdparagraph:P229:end

\end{minipage}\par\addvspace{4pt}

\par\begin{minipage}{\linewidth}\raggedright
% mdparagraph:P230:begin
\textbf{17.11.}
% mdparagraph:P230:end

% item: EX-IV-17-11
% key: KEY-EX-IV-17-11
% mdparagraph:P231:begin
He said to her, Give me, I beseech thee, water to drink; for I am thirsty.
% mdparagraph:P231:end

% mdparagraph:P232:begin
\textbf{Hē sægde hiere: Sęle mē, ic bidde ðē, wæter drincan; for-ðon-ðe ic eom ðurstig.}
% mdparagraph:P232:end

% mdparagraph:P233:begin
Hiere is the dative female addressee of sægde. Sęle is a singular command; mē is its dative recipient, while parenthetic ðē is the accusative person entreated by bidde. The visible help supplies \textbf{sęle mē wæter drincan}, with plain drincan and no tō. A grammatical placement \textbf{ic ðē bidde} is equivalent. Preserve the parenthesis and causal reason; do not treat every similar pronoun surface as having the same case.
% mdparagraph:P233:end

\end{minipage}\par\addvspace{4pt}

\par\begin{minipage}{\linewidth}\raggedright
% mdparagraph:P234:begin
\textbf{17.12.}
% mdparagraph:P234:end

% item: EX-IV-17-12
% key: KEY-EX-IV-17-12
% mdparagraph:P235:begin
He is not (om. a) God of (om. the) dead, but of (om. the) living.
% mdparagraph:P235:end

% mdparagraph:P236:begin
\textbf{Hē nis God dēadra, ac cwicra.}
% mdparagraph:P236:end

% mdparagraph:P237:begin
The supplied nis is the negative copula. Dēadra and cwicra are bare strong genitive plurals naming the dead and the living; the second shares the understood head God. Their articleless form preserves Cook's controls. Article-plus-weak forms could be grammatical in another phrasing, but would not meet those controls. The -ra forms are genitives, not nominative subjects.
% mdparagraph:P237:end

\end{minipage}\par\addvspace{4pt}

\par\begin{minipage}{\linewidth}\raggedright
% mdparagraph:P238:begin
\textbf{17.13.}
% mdparagraph:P238:end

% item: EX-IV-17-13
% key: KEY-EX-IV-17-13
% mdparagraph:P239:begin
The sun shines on (ofer) the righteous and the unrighteous.
% mdparagraph:P239:end

% mdparagraph:P240:begin
\textbf{Sēo sunne scīnð ofer ðā rihtwīsan and ðā unrihtwīsan.}
% mdparagraph:P240:end

% mdparagraph:P241:begin
Ofer governs both accusative plural personal-adjective groups. Their articles require weak rihtwīsan and unrihtwīsan, here naming righteous and unrighteous people. The second article may be shared in a defensible coordination. Keep the specified ofer rather than selecting a preposition solely from English on; bare strong endings would miss the selected definite construction.
% mdparagraph:P241:end

\end{minipage}\par\addvspace{4pt}

\par\begin{minipage}{\linewidth}\raggedright
% mdparagraph:P242:begin
\textbf{17.14.}
% mdparagraph:P242:end

% item: EX-IV-17-14
% key: KEY-EX-IV-17-14
% mdparagraph:P243:begin
Ye are (om. the) light of the world.
% mdparagraph:P243:end

% mdparagraph:P244:begin
\textbf{Gē sindon ðǣre worulde lēoht.}
% mdparagraph:P244:end

% mdparagraph:P245:begin
Gē has plural sindon. Lēoht is the bare nominative predicate, while ðǣre worulde is its genitive dependent. \textbf{Gē sindon lēoht ðǣre worulde} is also grammatical. Cook omits the article with light, not the world. The singular/\allowbreak{}mass predicate does not change the plural personal subject's agreement.
% mdparagraph:P245:end

\end{minipage}\par\addvspace{4pt}

\par\begin{minipage}{\linewidth}\raggedright
% mdparagraph:P246:begin
\textbf{17.15.}
% mdparagraph:P246:end

% item: EX-IV-17-15
% key: KEY-EX-IV-17-15
% mdparagraph:P247:begin
The twelve gates were twelve pearls.
% mdparagraph:P247:end

% mdparagraph:P248:begin
\textbf{Ðā twęlf gatu wǣron twęlf męregrēotan.}
% mdparagraph:P248:end

% mdparagraph:P249:begin
Twęlf is the familiar invariant numeral. Gatu is nominative plural subject; męregrēotan is nominative plural predicate identifying the gates as pearls, not a genitive partitive. Wǣron supplies the past. Another grammatical predicate order is acceptable. Do not replace gatu with singular geat or carry a genitive construction over from the higher numeral in 17.17.
% mdparagraph:P249:end

\end{minipage}\par\addvspace{4pt}

\par\begin{minipage}{\linewidth}\raggedright
% mdparagraph:P250:begin
\textbf{17.16.}
% mdparagraph:P250:end

% item: EX-IV-17-16
% key: KEY-EX-IV-17-16
% mdparagraph:P251:begin
The street of the city was gold.
% mdparagraph:P251:end

% mdparagraph:P252:begin
\textbf{Ðǣre ceastre strǣt wæs gold.}
% mdparagraph:P252:end

% mdparagraph:P253:begin
Strǣt is the singular nominative subject; ðǣre ceastre is its genitive dependent. Gold is a bare predicate noun. \textbf{Sēo strǣt ðǣre ceastre wæs gold} is a defensible alternative. The city is not the direct subject, and gold need not become an agreeing adjective.
% mdparagraph:P253:end

\end{minipage}\par\addvspace{4pt}

\par\begin{minipage}{\linewidth}\raggedright
% mdparagraph:P254:begin
\textbf{17.17.}
% mdparagraph:P254:end

% item: EX-IV-17-17
% key: KEY-EX-IV-17-17
% mdparagraph:P255:begin
Those that ate were four thousand men (154. c), besides women and children.
% mdparagraph:P255:end

% mdparagraph:P256:begin
\textbf{Ðā ðe ǣton wǣron fēower ðūsend manna, būtan wīfum and bearnum.}
% mdparagraph:P256:end

% mdparagraph:P257:begin
Ðā ðe ǣton identifies those who ate; both ǣton and wǣron are plural preterites. The printed help supplies the complete \textbf{fēower ðūsend manna} phrase, including genitive plural manna. Būtan governs dative wīfum and bearnum. Those who ate and the people who ate are equivalent interpretations. Do not replace the locally supplied genitive with an accusative or count the women and children inside the stated four thousand men.
% mdparagraph:P257:end

\end{minipage}\par\addvspace{4pt}

\reserve{7\baselineskip}
% mdparagraph:P258:begin
\section*{Cook, Group V, Exercise 18}
% mdparagraph:P258:end

\reserve{5\baselineskip}
\par\begin{minipage}{\linewidth}\raggedright
% mdparagraph:P259:begin
\textbf{18.1.}
% mdparagraph:P259:end

% item: EX-V-18-01
% key: KEY-EX-V-18-01
% mdparagraph:P260:begin
Blessed (p.p.) be God, the Father of our Lord Jesus Christ, the Father of mercies, and the God of all comfort.
% mdparagraph:P260:end

% mdparagraph:P261:begin
\textbf{Geblētsod sīe God, ūres Dryhtnes Hǣlendes Crīstes Fæder, mildheortnessa Fæder and ealre frōfre God.}
% mdparagraph:P261:end

% mdparagraph:P262:begin
Sīe is an independent blessing wish, with the required past participle geblētsod. God and the two Fæder phrases are nominative names of the same referent. Dryhtnes, Hǣlendes and Crīstes share genitive case in the title/\allowbreak{}name group; ūres belongs with Dryhtnes. Mildheortnessa is genitive plural mercies, and ealre agrees with feminine genitive singular frōfre. The genitives may follow their heads in a grammatical alternative. Indicative is would turn the wish into a statement, and nominal blētsung would miss Cook's participle control.
% mdparagraph:P262:end

\end{minipage}\par\addvspace{4pt}

\par\begin{minipage}{\linewidth}\raggedright
% mdparagraph:P263:begin
\textbf{18.2.}
% mdparagraph:P263:end

% item: EX-V-18-02
% key: KEY-EX-V-18-02
% mdparagraph:P264:begin
The winds were favorable.
% mdparagraph:P264:end

% mdparagraph:P265:begin
\textbf{Ðā windas wǣron gesyndige.}
% mdparagraph:P265:end

% mdparagraph:P266:begin
Windas is nominative plural with preterite wǣron. Strong masculine plural gesyndige means favorable here. The exact reading's syncopated gesyndge is a documented source variant; the full form is the selected production model. Singular wind or gesyndig would lose the number distinction, and healthy would miss the wind sense.
% mdparagraph:P266:end

\end{minipage}\par\addvspace{4pt}

\par\begin{minipage}{\linewidth}\raggedright
% mdparagraph:P267:begin
\textbf{18.3.}
% mdparagraph:P267:end

% item: EX-V-18-03
% key: KEY-EX-V-18-03
% mdparagraph:P268:begin
He wishes to watch their journey.
% mdparagraph:P268:end

% mdparagraph:P269:begin
\textbf{Hē wille hiera sīðfæt scēawian.}
% mdparagraph:P269:end

% mdparagraph:P270:begin
Wille expresses actual wishing and takes plain infinitive scēawian. Hiera is the plural possessor their; sīðfæt is the accusative journey watched. A defensible rearrangement is equivalent. Replacing hiera with his changes whose journey it is; scēawianne is the inflected infinitive and does not answer this willan-plus-plain-infinitive target.
% mdparagraph:P270:end

\end{minipage}\par\addvspace{4pt}

\par\begin{minipage}{\linewidth}\raggedright
% mdparagraph:P271:begin
\textbf{18.4.}
% mdparagraph:P271:end

% item: EX-V-18-04
% key: KEY-EX-V-18-04
% mdparagraph:P272:begin
The wind and the sea obey him.
% mdparagraph:P272:end

% mdparagraph:P273:begin
\textbf{Se wind and se sǣ hīerað him.}
% mdparagraph:P273:end

% mdparagraph:P274:begin
The two nominative subjects take plural hīerað. Sea uses the established masculine se sǣ. In the newly specified obey sense, him is dative. The verb and object may be ordered differently without changing the relation. Singular hīerð would ignore the coordination, while treating obey exactly like every English transitive verb would miss its dative government.
% mdparagraph:P274:end

\end{minipage}\par\addvspace{4pt}

\par\begin{minipage}{\linewidth}\raggedright
% mdparagraph:P275:begin
\textbf{18.5.}
% mdparagraph:P275:end

% item: EX-V-18-05
% key: KEY-EX-V-18-05
% mdparagraph:P276:begin
Not at all in vain fought I with death, as he fought with beasts.
% mdparagraph:P276:end

% mdparagraph:P277:begin
\textbf{Nālæs hōlunga feaht ic wið dēað, swā hē feaht wið dēor.}
% mdparagraph:P277:end

% mdparagraph:P278:begin
Nālæs hōlunga denies that the fighting was fruitless: not at all in vain. Both feaht forms are preterite indicative, first singular with ic and third singular with hē. Wið governs accusative opponents dēað and dēor; swā introduces the factual comparison as he fought. An equivalent grammatical order is acceptable. Do not add ne mechanically, translate the phrase as entirely in vain, or turn this as-clause into a purpose.
% mdparagraph:P278:end

\end{minipage}\par\addvspace{4pt}

\par\begin{minipage}{\linewidth}\raggedright
% mdparagraph:P279:begin
\textbf{18.6.}
% mdparagraph:P279:end

% item: EX-V-18-06
% key: KEY-EX-V-18-06
% mdparagraph:P280:begin
The time draws near [approaches, 164] (om. the) birth of Christ.
% mdparagraph:P280:end

% mdparagraph:P281:begin
\textbf{Sēo tīd genēalǣcð Crīstes gebyrde.}
% mdparagraph:P281:end

% mdparagraph:P282:begin
Genēalǣcð is third-singular present. The event approached is dative gebyrde, with genitive Crīstes; no article before birth is needed. A grammatical genitive-after-head order can preserve the meaning. Nominative gebyrd would miss genēalǣcan's dative complement, and making Christ the subject would reverse the stated relation.
% mdparagraph:P282:end

\end{minipage}\par\addvspace{4pt}

\par\begin{minipage}{\linewidth}\raggedright
% mdparagraph:P283:begin
\textbf{18.7.}
% mdparagraph:P283:end

% item: EX-V-18-07
% key: KEY-EX-V-18-07
% mdparagraph:P284:begin
It is the day when (ðe) he was born (ācenned).
% mdparagraph:P284:end

% mdparagraph:P285:begin
\textbf{Hit is se dæg ðe hē wæs ācenned.}
% mdparagraph:P285:end

% mdparagraph:P286:begin
Se dæg is the antecedent of the temporal relative ðe; hē wæs ācenned says he was born. The exact supplied participle retains plain e. The day on which he was born gives the same interpretation. Interrogative hwonne would not satisfy the relative target, and the supplied participle is not evidence for unaided production of a complete birth-verb paradigm.
% mdparagraph:P286:end

\end{minipage}\par\addvspace{4pt}

\par\begin{minipage}{\linewidth}\raggedright
% mdparagraph:P287:begin
\textbf{18.8.}
% mdparagraph:P287:end

% item: EX-V-18-08
% key: KEY-EX-V-18-08
% mdparagraph:P288:begin
O, wast thou with me, dearest, then, while I rose up against my doom (dōm)?
% mdparagraph:P288:end

% mdparagraph:P289:begin
\textbf{Ēalā, wǣre ðū mid mē, lēofosta, ðā, ðā-hwīle-ðe ic rās ūp ongēan mīnne dōm?}
% mdparagraph:P289:end

% mdparagraph:P290:begin
Wǣre is second-singular preterite indicative in this factual past question; its shape alone does not make it subjunctive. Mid mē expresses accompaniment. Lēofosta is weak masculine singular superlative in direct address; feminine \textbf{lēofoste} is equally defensible for an unstated female addressee. Retain then, while, up and against: ðā, ðā-hwīle-ðe, ūp and ongēan mīnne dōm. Rās is first-singular preterite, and the opposed doom/\allowbreak{}fate is accusative. Do not mistake dearest for a comparative or omit the rising up.
% mdparagraph:P290:end

\end{minipage}\par\addvspace{4pt}

\par\begin{minipage}{\linewidth}\raggedright
% mdparagraph:P291:begin
\textbf{18.9.}
% mdparagraph:P291:end

% item: EX-V-18-09
% key: KEY-EX-V-18-09
% mdparagraph:P292:begin
The moon is hid (hȳdan, 113), the night is still.
% mdparagraph:P292:end

% mdparagraph:P293:begin
\textbf{Se mōna is hȳded, sēo niht is stillu.}
% mdparagraph:P293:end

% mdparagraph:P294:begin
Mōna is a singular masculine subject with past-participle hȳded. The documented contracted participle hȳd is also valid here. Niht is feminine nominative singular with strong positive stillu. These are two predicates of different kinds: hidden is participial, still is an ordinary adjective. Dictionary headword stille does not automatically supply the requested feminine predicate form.
% mdparagraph:P294:end

\end{minipage}\par\addvspace{4pt}

\par\begin{minipage}{\linewidth}\raggedright
% mdparagraph:P295:begin
\textbf{18.10.}
% mdparagraph:P295:end

% item: EX-V-18-10
% key: KEY-EX-V-18-10
% mdparagraph:P296:begin
Thou bringst the fisher to his wife.
% mdparagraph:P296:end

% mdparagraph:P297:begin
\textbf{Ðū bringst ðone fiscere tō his wīfe.}
% mdparagraph:P297:end

% mdparagraph:P298:begin
Bringst agrees with singular ðū. The brought person is accusative ðone fiscere; tō takes dative wīfe as the destination/\allowbreak{}person. His identifies the fisherman as the wife's possessor. A grammatical order such as \textbf{Ðū bringst tō his wīfe ðone fiscere} preserves the relation. Do not treat the fisherman as dative merely because the destination is dative too.
% mdparagraph:P298:end

\end{minipage}\par\addvspace{4pt}

\par\begin{minipage}{\linewidth}\raggedright
% mdparagraph:P299:begin
\textbf{18.11.}
% mdparagraph:P299:end

% item: EX-V-18-11
% key: KEY-EX-V-18-11
% mdparagraph:P300:begin
Calm is the morn(ing).
% mdparagraph:P300:end

% mdparagraph:P301:begin
\textbf{Smylte is se morgen.}
% mdparagraph:P301:end

% mdparagraph:P302:begin
Smylte is the strong masculine nominative singular predicate agreeing with se morgen. \textbf{Se morgen is smylte} preserves the meaning in subject-first order. The initial adjective is not an adverb modifying is; an adverbial form would miss the predicate-agreement target.
% mdparagraph:P302:end

\end{minipage}\par\addvspace{4pt}

\par\begin{minipage}{\linewidth}\raggedright
% mdparagraph:P303:begin
\textbf{18.12.}
% mdparagraph:P303:end

% item: EX-V-18-12
% key: KEY-EX-V-18-12
% mdparagraph:P304:begin
I know that this was life.
% mdparagraph:P304:end

% mdparagraph:P305:begin
\textbf{Ic wāt ðæt ðis wæs līf.}
% mdparagraph:P305:end

% mdparagraph:P306:begin
Wāt is first-singular present: the knowing is current, while the content \textbf{ðæt ðis wæs līf} identifies something in the past. Ðis is neuter nominative singular and līf the predicate noun. A grammatical change in clause order is acceptable if the present knowing and past identification remain. Life is not light; making ðæt a purpose conjunction or replacing wæs with a future form changes the statement.
% mdparagraph:P306:end

\end{minipage}\par\addvspace{4pt}

\par\begin{minipage}{\linewidth}\raggedright
% mdparagraph:P307:begin
\textbf{18.13.}
% mdparagraph:P307:end

% item: EX-V-18-13
% key: KEY-EX-V-18-13
% mdparagraph:P308:begin
What do ye more than they?
% mdparagraph:P308:end

% mdparagraph:P309:begin
\textbf{Hwæt dōð gē mā ðonne hīe?}
% mdparagraph:P309:end

% mdparagraph:P310:begin
Hwæt is the accusative object of second-plural dōð. Mā compares the amount of doing; ðonne hīe has nominative hīe with the repeated do understood. \textbf{Hwæt dōð gē mā ðonne hīe dōð?} makes that predicate explicit. Do not replace adverbial mā with an adjective agreeing with a noun, or make the comparison standard accusative merely because hwæt is the object.
% mdparagraph:P310:end

\end{minipage}\par\addvspace{4pt}

\reserve{7\baselineskip}
% mdparagraph:P311:begin
\section*{R. The fisherman and his son}
% mdparagraph:P311:end

\reserve{5\baselineskip}
\par\begin{minipage}{\linewidth}\raggedright
% mdparagraph:P312:begin
\textbf{R01. Complete translation.}
% mdparagraph:P312:end

% item: NEW-U11-R01
% key: KEY-NEW-U11-R01
% mdparagraph:P313:begin
The old fisherman came to the house so that he might help his son. When he came, the son was hungry. The fisherman gave him food. While the son ate, the fisherman carried water. Although the fisherman was old, he helped his son. The son said, “Be here until I come.” He went away and came back and gave his father water.
% mdparagraph:P313:end

% mdparagraph:P314:begin
\textbf{Interpretation:} Preserve the purpose of the first arrival, the simultaneous eating and carrying, the concession about age, and the son's answering act of help. In order to help his son, as the son ate and stay here are equivalent English choices. Came back is a contextual rendering of cōm after the departure, not a separate return prefix. The final he is most naturally the son in this sequence; masculine gender alone does not rule out another referent. Do not omit the final gift or turn the admitted age into a condition.
% mdparagraph:P314:end

\end{minipage}\par\addvspace{4pt}

\par\begin{minipage}{\linewidth}\raggedright
% mdparagraph:P315:begin
\textbf{R02. Three relations in the account.}
% mdparagraph:P315:end

% item: NEW-U11-R02
% key: KEY-NEW-U11-R02
% mdtable:T13:begin
\par\begingroup
\fontsize{10.5}{13.2}\selectfont
\setlength{\LTleft}{0pt}\setlength{\LTright}{\fill}
\noindent
\begin{tabular}{@{}>{\raggedright\arraybackslash}p{21.5mm}>{\raggedright\arraybackslash}p{22.78mm}>{\raggedright\arraybackslash}p{74.09mm}@{}}
\toprule
\textbf{Relation} & \textbf{Clause} & \textbf{Contribution and finite form} \\
\midrule
Purpose & ðæt hē his suna hulpe & Why the fisherman came; hulpe is third-singular preterite subjunctive. \\
Simultaneous time & Mid-ðȳ se sunu ǣt & The son eats while his father carries water; ǣt is third-singular preterite indicative. \\
Concession & Ðēah se fiscere wǣre eald & Age is admitted despite which the help occurs; wǣre is third-singular preterite subjunctive. \\
\bottomrule
\end{tabular}
\par\endgroup

% mdtable:T13:end
% mdparagraph:P316:begin
Purpose/\allowbreak{}intention, while/\allowbreak{}as and although/\allowbreak{}though are equivalent labels or translations when these distinctions remain. Calling all three clauses conditions loses their separate functions. The concessive subjunctive does not make the father's age false.
% mdparagraph:P316:end

\end{minipage}\par\addvspace{4pt}

\par\begin{minipage}{\linewidth}\raggedright
% mdparagraph:P317:begin
\textbf{R03. Givers, recipients and possessors.}
% mdparagraph:P317:end

% item: NEW-U11-R03
% key: KEY-NEW-U11-R03
% mdparagraph:P318:begin
The fisherman gives the first food, and his son receives it. \textbf{Se fiscere} explicitly names that giver; \textbf{him} is the dative recipient understood as the son. In \textbf{his suna}, his refers to the fisherman and suna is the dative son helped.
% mdparagraph:P318:end

% mdparagraph:P319:begin
At the end, the son is the natural subject of \textbf{hē … sealde}, and his father receives the water. \textbf{His fæder} is a dative phrase; fæder has the selected endingless dative singular, while his is the son's genitive possessor. The pronouns' forms alone do not identify one of the two men uniquely. A reasoned alternative antecedent analysis should be distinguished from the account's intended reciprocal sequence, rather than rejected as morphologically impossible. Treating endingless fæder as necessarily nominative would make the recipient a new subject without justification.
% mdparagraph:P319:end

\end{minipage}\par\addvspace{4pt}

\par\begin{minipage}{\linewidth}\raggedright
% mdparagraph:P320:begin
\textbf{R04. The son’s expected return.}
% mdparagraph:P320:end

% item: NEW-U11-R04
% key: KEY-NEW-U11-R04
% mdparagraph:P321:begin
\textbf{Ic} is the son, whose speech is introduced by \textbf{se sunu cwæð}. When he tells his father \textbf{wes hēr}, his own coming is still expected. \textbf{Cume} is first-singular present subjunctive in the selected anticipated-until clause, formally identical to first-singular indicative cume. The later \textbf{cōm}, after \textbf{gewāt}, fulfills that expectation on the natural reading of the final hē as the son.
% mdparagraph:P321:end

% mdparagraph:P322:begin
Until I arrive is an equivalent rendering. The first person belongs to the quoted speaker, not to the narrator or father; a claim that the ending of cume by itself proves the mood would overstate the morphology.
% mdparagraph:P322:end

\end{minipage}\par\addvspace{4pt}

\reserve{7\baselineskip}
% mdparagraph:P323:begin
\section*{B. Britain: metals and jet}
% mdparagraph:P323:end

\reserve{5\baselineskip}
\par\begin{minipage}{\linewidth}\raggedright
% mdparagraph:P324:begin
\textbf{B01. Complete translation.}
% mdparagraph:P324:end

% item: NEW-U11-B01
% key: KEY-NEW-U11-B01
% mdparagraph:P325:begin
Likewise it is also rich in veins of metal—copper and iron, lead and silver. Jet is also found here; the stone is a black gem. If one puts it in fire, then snakes flee away from there.
% mdparagraph:P325:end

% mdparagraph:P326:begin
\textbf{Interpretation:} Initial it is Britain, identified in the introduction. Productive in ores of metals is a defensible rendering of the supported first expression; the four named metals must remain. Gagates is jet, not the agent doing the finding. Generic man may be rendered one, someone or a person. Adders is a valid alternative to snakes. Flēoð means flee, not fly; the conditional claim must not become an unconditional report that someone put the stone in the fire. [5]
% mdparagraph:P326:end

\end{minipage}\par\addvspace{4pt}

\par\begin{minipage}{\linewidth}\raggedright
% mdparagraph:P327:begin
\textbf{B02. The condition and its consequence.}
% mdparagraph:P327:end

% item: NEW-U11-B02
% key: KEY-NEW-U11-B02
% mdparagraph:P328:begin
The condition is \textbf{gif man hine on fȳr dēð}, if one puts it in fire. Its main clause is \textbf{ðonne flēoð ðǣr næddran onwęg}, then snakes flee away from there. \textbf{Dēð} is third-singular present indicative of dōn; \textbf{flēoð} is third-plural present indicative of flēon. \textbf{Ðonne} is the main-clause then, correlating with the condition.
% mdparagraph:P328:end

% mdparagraph:P329:begin
When someone puts it in fire can express the same general relationship idiomatically, but the requested grammatical label is the gif-condition. Ðonne is neither comparative than nor the subject of flēoð. The plural snakes control that verb; generic man controls singular dēð.
% mdparagraph:P329:end

\end{minipage}\par\addvspace{4pt}

\par\begin{minipage}{\linewidth}\raggedright
% mdparagraph:P330:begin
\textbf{B03. Two different referents.}
% mdparagraph:P330:end

% item: NEW-U11-B03
% key: KEY-NEW-U11-B03
% mdparagraph:P331:begin
Initial \textbf{hit} refers to Britain, the island identified in the introduction: it is neuter nominative singular in the copular clause. \textbf{Hine} refers to \textbf{se stān}, the masculine stone just named, and is accusative singular as the thing put into the fire. Thus the passage changes referent.
% mdparagraph:P331:end

% mdparagraph:P332:begin
Britain/\allowbreak{}the island and the stone/\allowbreak{}the piece of jet are equivalent contextual descriptions. English it can conceal the gender and case contrast; that does not make the two Old English pronouns interchangeable or make both refer to the island.
% mdparagraph:P332:end

\end{minipage}\par\addvspace{4pt}

\clearpage
% mdparagraph:P333:begin
\section*{S. Æthelwald calms the sea}
% mdparagraph:P333:end

\reserve{5\baselineskip}
\par\begin{minipage}{\linewidth}\raggedright
% mdparagraph:P334:begin
\textbf{S01. Complete translation.}
% mdparagraph:P334:end

% item: NEW-U11-S01
% key: KEY-NEW-U11-S01
% mdparagraph:P335:begin
Then, after a long time, when we raised our gaze from a distance, we saw on Farne, the island, the father Æthelwald, dear to God, coming out of his secluded place, so that he might watch our journey and see what might happen to us, for he heard the crashing of the storms and of the seething sea.
% mdparagraph:P335:end

% mdparagraph:P336:begin
When he then also looked at us and saw us placed in distress and despair, he bent his knees to the Father of our Lord Jesus Christ and prayed for our safety and our life. And when he completed that prayer, he then at the same time both calmed the swollen sea and stilled the storm, so that the violence of the storm ceased entirely, and favorable winds brought us to land through the very calm sea. When we then came up onto land and also carried our ship up from the waves, the same storm soon returned and came again; it had ceased for a short time for our sake, and it was very great and strong all that day. Thus people could clearly understand that the short period of calm which had occurred there was granted from heaven at the prayers of that man of God for our rescue.
% mdparagraph:P336:end

% mdparagraph:P337:begin
\textbf{Interpretation:} The opening can be recast as After a long time we lifted our gaze and saw …, provided the distant seeing and temporal sequence remain. Æthelwald is the person seen; \textbf{Gode} belongs with \textbf{lēofan}, dear to God. Our Lord and Saviour Christ is an alternative rendering of the title group. Our lives is idiomatic English for the collective situation, although līfe is morphologically singular.
% mdparagraph:P337:end

% mdparagraph:P338:begin
The subject \textbf{hē} in the paired calming and stilling actions remains Æthelwald. The narrative attributes efficacy to his prayer; a translation should not silently replace the grammatical subject with God. Prayed and ceased are valid simple-past renderings of \textbf{wæs gebiddende} and \textbf{wæs blinnende}, as Cook's notes permit. A contextual English progressive is possible, but is not forced by the participle's form. Keep the landing, ship carried ashore, storm's return and final inference. Final \textbf{forgiefen} means given/\allowbreak{}granted, not forgiven. [5]
% mdparagraph:P338:end

\end{minipage}\par\addvspace{4pt}

\par\begin{minipage}{\linewidth}\raggedright
% mdparagraph:P339:begin
\textbf{S02. The order of the rescue.}
% mdparagraph:P339:end

% item: NEW-U11-S02
% key: KEY-NEW-U11-S02
\begin{itemize}
% mdparagraph:P340:begin
\item The sailors see Æthelwald.
% mdparagraph:P340:end
% mdparagraph:P341:begin
\item Æthelwald prays.
% mdparagraph:P341:end
% mdparagraph:P342:begin
\item The sea becomes calm.
% mdparagraph:P342:end
% mdparagraph:P343:begin
\item The sailors reach land.
% mdparagraph:P343:end
% mdparagraph:P344:begin
\item The storm returns.
% mdparagraph:P344:end
\end{itemize}

% mdparagraph:P345:begin
Before its return they have both come ashore and carried their ship up from the waves. Reach land/\allowbreak{}come ashore are equivalent descriptions, but the ship's removal is a separate completed action in the same temporal clause. Putting the renewed storm before the landing, or stopping the episode at the calm, loses the resolution.
% mdparagraph:P345:end

\end{minipage}\par\addvspace{4pt}

\par\begin{minipage}{\linewidth}\raggedright
% mdparagraph:P346:begin
\textbf{S03. Purpose, question and reason.}
% mdparagraph:P346:end

% item: NEW-U11-S03
% key: KEY-NEW-U11-S03
% mdparagraph:P347:begin
The purpose begins \textbf{ðæt hē wolde ūrne sīðfæt scēawian, and gesēon hwæt ūs gelumpe}: so that he might watch our journey and see what might happen to us. The local help supplies wolde with its two infinitives in this purpose use; in order to watch and see is also an apt translation. \textbf{Hwæt ūs gelumpe} is the dependent question inside that purpose, not a second outer purpose clause.
% mdparagraph:P347:end

% mdparagraph:P348:begin
The reason begins \textbf{for-ðon hē gehīerde ðæt gebrec ðǣra storma and ðæs weallendes sǣs}: because he heard the crashing of the storms and seething sea. It states the circumstance prompting him to come out. A rendering with since is valid if clearly causal. Reducing wolde to private desire alone, or treating the heard noise as an intended outcome, would lose the contrast.
% mdparagraph:P348:end

\end{minipage}\par\addvspace{4pt}

\par\begin{minipage}{\linewidth}\raggedright
% mdparagraph:P349:begin
\textbf{S04. The narrated result.}
% mdparagraph:P349:end

% item: NEW-U11-S04
% key: KEY-NEW-U11-S04
% mdparagraph:P350:begin
\textbf{Tō ðon ðætte} introduces the result: the storm's violence ceased entirely, and favorable winds brought the sailors to land through the very calm sea. The sailors are \textbf{ūs}, the accusative people brought; \textbf{gesyndge windas} is the subject of \textbf{gebrōhton}.
% mdparagraph:P350:end

% mdparagraph:P351:begin
The supplied \textbf{wæs blinnende} can be translated ceased. Was ceasing may be defensible English in context, but no progressive translation is required merely by the present participle. The two coordinated outcomes are narrated as results, not just things Æthelwald hoped would happen. Leaving out the winds would obscure the expressed agent of the bringing.
% mdparagraph:P351:end

\end{minipage}\par\addvspace{4pt}

\par\begin{minipage}{\linewidth}\raggedright
% mdparagraph:P352:begin
\textbf{S05. Where the temporal clauses end.}
% mdparagraph:P352:end

% item: NEW-U11-S05
% key: KEY-NEW-U11-S05
% mdparagraph:P353:begin
The first requested clause extends from \textbf{Mid-ðȳ hē ðā ūs ēac scēawode} through \textbf{in gewinne and in ōrmōdnesse gesętte bēon}. It includes both his looking at them and his seeing their plight. The main action begins \textbf{ðā biegde hē his cnēowu}, then he bent his knees.
% mdparagraph:P353:end

% mdparagraph:P354:begin
The second extends from \textbf{Mid-ðȳ wē ðā ūp cōmon tō lande} through \textbf{fram ðǣm ȳðum ūp ābǣron}. It includes coming ashore and carrying the ship up. The main action begins \textbf{ðā sōna se ilca storm ęft hwearf and cōm}, then the same storm soon returned and came again.
% mdparagraph:P354:end

% mdparagraph:P355:begin
When or while can represent mid-ðȳ as context permits. The relevant boundaries, rather than one fixed English conjunction, are the target. Ending the second subordinate clause at lande wrongly detaches the ship's removal from the circumstances completed before the storm returns.
% mdparagraph:P355:end

\end{minipage}\par\addvspace{4pt}

\par\begin{minipage}{\linewidth}\raggedright
% mdparagraph:P356:begin
\textbf{S06. Wæs gebiddende.}
% mdparagraph:P356:end

% item: NEW-U11-S06
% key: KEY-NEW-U11-S06
% mdparagraph:P357:begin
Here \textbf{wæs gebiddende} can be translated \textbf{prayed}. Cook's note points back to the explanation that this be-plus-present-participle construction can have simple-past force. \textbf{Was praying} is an idiomatic contextual alternative, but the present participle does not require the English progressive category.
% mdparagraph:P357:end

% mdparagraph:P358:begin
The auxiliary wæs gives the past setting; calling the whole event present because gebiddende is a present participle would confuse nonfinite form with finite tense. This local explanation supports the source passage without making every Old English be-plus-participle phrase equivalent to an English progressive. [5]
% mdparagraph:P358:end

\end{minipage}\par\addvspace{4pt}

\par\begin{minipage}{\linewidth}\raggedright
% mdparagraph:P359:begin
\textbf{S07. The meaning of the final inference.}
% mdparagraph:P359:end

% item: NEW-U11-S07
% key: KEY-NEW-U11-S07
% mdparagraph:P360:begin
People understand that the brief calm was granted from heaven, at the prayers of the man of God, for the sailors' rescue. \textbf{Ðæs Godes weres} identifies whose prayers are meant: Godes depends on weres, not on bēnum as though God were the person praying. \textbf{Ūre} in \textbf{for intingan ūre hǣle} identifies the sailors whose rescue is at issue.
% mdparagraph:P360:end

% mdparagraph:P361:begin
Given/\allowbreak{}granted from heaven and for our safety/\allowbreak{}rescue are defensible translations here. Cook explicitly explains final \textbf{forgiefen} as given; forgiven would introduce a meaning that this clause does not express. The content of the inference concerns the specially granted interval of calm, not a permanent end to all storms.
% mdparagraph:P361:end

\end{minipage}\par\addvspace{4pt}

\par\begin{minipage}{\linewidth}\raggedright
% mdparagraph:P362:begin
\textbf{S08. Why the storm returns in the account.}
% mdparagraph:P362:end

% item: NEW-U11-S08
% key: KEY-NEW-U11-S08
% mdparagraph:P363:begin
The calm is temporary. After the sailors and their ship are ashore, the same storm returns and remains very great and strong throughout that day. Its short suspension precisely when they needed rescue supports the narrator's explanation that the calm was granted in response to prayer for their safety.
% mdparagraph:P363:end

% mdparagraph:P364:begin
A paraphrase such as the storm paused long enough for their escape preserves the point. Saying that the storm ended permanently conflicts with its explicit return; saying that the narrator sees no relation to the prayer omits the final inference.
% mdparagraph:P364:end

\end{minipage}\par\addvspace{4pt}

\clearpage
% mdparagraph:P365:begin
\section*{Source notes}
% mdparagraph:P365:end

\begingroup\fontsize{10}{12.8}\selectfont
% mdparagraph:P366:begin
[1] Cook, \emph{Exercises in Old English} (1895), Group IV, Exercises 16–17, and Group V, Exercise 18, printed pp. 8–9; one-based PDF pp. 16–17. Original prompts are retained. All answers and explanatory prose are new teaching material.
% mdparagraph:P366:end

% mdparagraph:P367:begin
[2] Cook, \emph{A First Book in Old English}, second edition (1895), §§194–196, printed pp. 103–105; Fulk, \emph{An Introductory Grammar of Old English} (2020), §§44–45, printed pp. 28–30. The selected purpose/\allowbreak{}result distinction and factual concessive use follow the reconciled primary evidence, rather than a rule assigning one mood to every occurrence of a conjunction.
% mdparagraph:P367:end

% mdparagraph:P368:begin
[3] Cook, \emph{First Book}, the noun and adjective patterns on printed pp. 28–43, especially §43.4–5 on the season nouns and §58.2 on consonant simplification, with the vocabulary entries cited in the learner's notes. The unassigned continuation on p. 139 attests geatum; it is a valid variant, not an additional reading assignment.
% mdparagraph:P368:end

% mdparagraph:P369:begin
[4] Cook, \emph{First Book}, the selected strong and weak verb systems and participles, printed pp. 54–78; the lexical entries and documented alternatives are specified in the learner's source notes. Complete requested cells are printed above; valid variants do not establish an additional conjugational system.
% mdparagraph:P369:end

% mdparagraph:P370:begin
[5] Cook, \emph{First Book}, p. 139, editorial prose lines 1–4, and pp. 142–143, from editorial line 8 of p. 142 through editorial line 14 of p. 143; one-based PDF pp. 155 and 158–159. On p. 143, note 1 is attached to both wæs gebiddende and wæs blinnende and refers back to p. 142 note 7; p. 143 note 18 explains forgiefen as given rather than forgiven. The complete translations above were independently prepared from the Old English before comparison with the other private witnesses.
% mdparagraph:P370:end

\endgroup
\end{document}
